% ===========================================================================
% Chapitre 02 — React vu par l'analyste : le modèle d'exécution que l'on abstrait
% Sources : notes/16-react-semantics.md ; ADR-001, 004, 009, 017, 025, 026, 035.
% Sorties d'analyseur observées au commit e67b10a (cargo run -q -- check ...),
% fichiers d'exemple sous /tmp/ch02/ et /tmp/ch02next/.
% ===========================================================================
\chapter{React vu par l'analyste : le modèle d'exécution que l'on abstrait}
\label{chap:react-modele}

\begin{objectifs}
  \item Décrire précisément ce que React fait d'un composant : déclenchement,
    rendu, commit, effets, et la façon dont un \code{setState} relance le cycle.
  \item Connaître les faits de React sur lesquels repose toute l'analyse :
    file d'updaters et batching, comparaison \code{Object.is} des deps et de
    l'état, identité fraîche des littéraux, ordre positionnel des hooks,
    Strict Mode, contexte, Server Components.
  \item Lire la sémantique formelle React-tRace (syntaxe, mémoire d'arbre,
    transitions \code{StepInit}/\code{StepEffect}/\code{StepCheck}/\code{StepEvent},
    théorèmes) et savoir pourquoi le projet l'a prise comme sémantique concrète
    de référence (\adr{1}).
  \item Savoir, pour chaque fait de React, ce que \reactant{} en garde, ce qu'il
    en abstrait, où cela se trouve dans le code, et pourquoi l'abstraction
    reste une sur-approximation.
  \item Connaître les exemples canoniques (\code{Counter}, \code{SelfCounter},
    \code{Inf}, \code{Timer}, \code{Toggle}, \code{ObjChurn}\dots) réutilisés dans
    tout le manuscrit, avec la sortie réelle de l'analyseur sur chacun.
\end{objectifs}

\prerequis{\cref{chap:introduction} (les trois programmes fautifs et la promesse
\Error/\Warning/\Info). Aucune connaissance d'analyse statique n'est requise : les
notions de treillis, de join et de soundness sont seulement nommées ici et
définies au \cref{chap:interpretation-abstraite}.}

Un analyseur par interprétation abstraite ne calcule pas « dans le vide » : il
sur-approxime une \terme{sémantique concrète}, c'est-à-dire une description
exacte de ce que fait le programme lorsqu'il s'exécute. Pour un programme React,
« ce que fait le programme » n'est pas seulement ce que fait le code JavaScript
du composant : c'est ce que fait React \emph{avec} ce code --- quand il appelle la
fonction, ce qu'il garde entre deux appels, quand il exécute les effets, ce qu'il
compare pour décider de recommencer. Ce chapitre décrit ce modèle d'exécution du
point de vue de quelqu'un qui doit l'abstraire.

Trois sources sont distinguées tout au long du chapitre, et il importe de ne
jamais les confondre :
\begin{itemize}
  \item \textbf{React} --- la documentation officielle \cite{ReactDocs}, qui fixe
    le comportement attendu (les citations sont reprises en anglais, telles
    quelles) ;
  \item \textbf{React-tRace} --- l'article de Lee, Ahn et Yi (OOPSLA 2025)
    \cite{LeeAhnYi2025}, qui donne une sémantique opérationnelle formelle d'un
    noyau de React ;
  \item \textbf{le code} --- ce que \reactant{} calcule réellement au commit
    \code{e67b10a}, localisé par fichier et lignes.
\end{itemize}
Les huit premières sections suivent la progression d'un utilisateur de React :
rendu, état, effets, identité, règles des hooks, Strict Mode, contexte, Server
Components. Chacune se termine par ce que l'analyseur en fait. Les trois
dernières prennent de la hauteur : la sémantique formelle
(\cref{sec:react-modele-trace}), la sémantique concrète que le projet adopte
réellement (\cref{sec:react-modele-concrete}) et le catalogue des exemples du
manuscrit (\cref{sec:react-modele-exemples}).

\section{Un composant est une fonction du rendu}
\label{sec:react-modele-rendu}

\subsection{Un premier programme : \code{Counter}}

Commençons par le programme React le plus simple qui ait un état. C'est le
premier exemple de l'article React-tRace, transcrit en TSX.

\begin{exemple}{\code{Counter} : un état, un handler}{canon-counter}
\begin{lstlisting}[language=TypeScript,caption={\code{Counter} (fichier \file{/tmp/ch02/counter.tsx})},label={lst:react-modele-counter}]
import { useState } from "react";

export function Counter() {
  const [count, setCount] = useState(0);
  return (
    <button onClick={() => setCount((c) => c + 1)}>
      {count}
    </button>
  );
}
\end{lstlisting}
\end{exemple}

Déroulons ce qui se passe à l'exécution.
\begin{enumerate}
  \item \emph{Montage.} React appelle la fonction \code{Counter()}. L'appel
    \code{useState(0)} ne trouve aucun état existant : React en crée un, lui
    donne la valeur initiale \code{0}, et renvoie la paire
    \code{[0, setCount]}. La fonction renvoie un élément JSX, c'est-à-dire une
    \emph{description} d'un bouton dont le texte vaut \code{0}.
  \item \emph{Commit.} React compare cette description à ce qui est affiché
    (rien) et crée le nœud DOM \code{<button>}.
  \item \emph{Clic.} Le navigateur déclenche le handler
    \code{() => setCount((c) => c + 1)}. L'appel à \code{setCount} ne modifie
    pas \code{count} : il \emph{met en file} une mise à jour et demande un
    nouveau rendu.
  \item \emph{Second rendu.} React rappelle \code{Counter()}. Cette fois
    \code{useState(0)} retrouve l'état existant, applique la mise à jour en
    file (\code{c => c + 1} appliqué à \code{0}) et renvoie \code{[1, setCount]}.
    L'argument \code{0} est ignoré.
  \item \emph{Second commit.} React compare la nouvelle description à
    l'ancienne ; seul le texte du bouton diffère, seul ce nœud texte est
    modifié.
\end{enumerate}

\begin{react}[déclencher, rendre, valider]
La documentation décrit trois étapes \cite{ReactRenderCommit} :
\og Triggering a render\fg{}, \og Rendering the component\fg{} et
\og Committing to the DOM\fg{}. Il n'y a que deux raisons de rendre un
composant : \og It's the component's initial render\fg{} et \og The component's
(or one of its ancestors') state has been updated\fg{}. Le rendu est récursif :
rendre un composant rend les composants qu'il renvoie. Enfin, \og React only
changes the DOM nodes if there's a difference between renders\fg{}.
\end{react}

Pour l'analyste, la conséquence est simple et structurante : \emph{le composant
est une fonction que React appelle autant de fois qu'il le veut}. Entre deux
appels, rien ne survit dans les variables locales de la fonction ; tout ce qui
survit est rangé par React, à côté du composant, et rendu à l'appel suivant par
les hooks.

\begin{definition}{Tour de React}{tour-react}
Un \terme{tour} de React est la séquence
\[
  \text{déclencheur} \;\to\; \Render \;\to\; \Commit \;\to\; \Effect
\]
où le \emph{déclencheur} est le montage ou une mise à jour d'état ; le
\emph{rendu} ($\Render$) appelle les fonctions des composants concernés et
produit une description de l'interface ; le \emph{commit} ($\Commit$) applique
au DOM la différence avec la description précédente ; les \emph{effets}
($\Effect$) exécutent les callbacks \code{useEffect} dont les deps ont changé.
Si les effets appellent un setter, un nouveau tour est déclenché. Entre deux
tours, React attend un événement extérieur (clic, réponse réseau, minuterie) :
c'est la \emph{boucle d'événements}.
\end{definition}

La \cref{fig:react-modele-chronologie} représente un tour et les deux façons
d'en déclencher un autre. Elle servira de référence dans tout le manuscrit :
les phases $\Render$, $\Commit$, $\Effect$ et le temps « plus tard » ($\Tick$,
un tour ultérieur de la boucle d'événements) sont celles que la relation des
écrivains distingue (\cref{chap:relations}).

\ifhastikz
\begin{figure}[htbp]\centering
\begin{tikzpicture}[node distance=7mm and 5mm]
  \node[bloc] (trig) {déclencheur\\\scriptsize montage ou \code{set}};
  \node[bloc,right=of trig] (render) {rendu\\\scriptsize appel du corps};
  \node[bloc,right=of render] (commit) {commit\\\scriptsize DOM modifié};
  \node[bloc,right=of commit] (paint) {peinture\\\scriptsize navigateur};
  \node[bloc,right=of paint] (eff) {effets passifs\\\scriptsize \code{useEffect}};
  \draw[fleche] (trig) -- (render);
  \draw[fleche] (render) -- (commit);
  \draw[fleche] (commit) -- (paint);
  \draw[fleche] (paint) -- (eff);
  \path (commit.north) -- (paint.north) coordinate[midway] (mcp);
  \path (render.south) -- (paint.south) coordinate[midway] (mrp);
  \node[etiquette,above=1mm of mcp] {\code{useLayoutEffect}};
  \draw[flechemust] ([xshift=-4mm]eff.south) -- ++(0,-0.9)
     -| node[etiquette,pos=0.25,above]{\code{setState} dans un effet : nouveau tour} (trig.south);
  \node[bloc,below=18mm of mrp] (loop) {boucle d'événements\\\scriptsize attente d'un clic, d'un timer, d'une réponse};
  \draw[fleche] ([xshift=4mm]eff.south) |- node[etiquette,pos=0.3,right]{aucun \code{set}} (loop.east);
  \draw[flechemay] (loop.west) -| node[etiquette,pos=0.3,below]{handler ou timer : \code{set}} ([xshift=-5mm]trig.south);
  \draw[flechemust] ([xshift=-4mm]render.north) .. controls +(-0.4,1.1) and +(0.4,1.1) ..
     node[etiquette,above]{\code{set} pendant le rendu : rendu jeté, on recommence} ([xshift=4mm]render.north);
\end{tikzpicture}
\caption{Chronologie d'un tour de React. Un setter appelé dans un effet relance
immédiatement un tour (flèche pleine) : c'est le mécanisme de la boucle
infinie d'effets. Un setter appelé dans un handler ne relance un tour qu'après
un événement extérieur (flèche tiretée). Un setter appelé pendant le rendu fait
jeter le rendu et recommencer.}
\label{fig:react-modele-chronologie}
\end{figure}
\fi

\subsection{Le rendu doit être pur}

\begin{react}[pureté du rendu]
\og Rendering must always be a pure calculation\fg{} \cite{ReactRenderCommit} :
mêmes entrées, même sortie (\og Same inputs, same output\fg{}), et le composant
ne modifie rien qui existait avant son appel (\og It minds its own business\fg{}).
Les entrées d'un rendu sont les props, l'état et le contexte.
\end{react}

La pureté est ce qui rend l'analyse possible : si le corps d'un composant est
une fonction de (props, état, contexte), l'analyseur peut l'évaluer
abstraitement sur des \emph{ensembles} de props et d'états sans se soucier du
nombre de fois où React l'appelle. Mais la pureté n'est qu'une convention ---
rien n'empêche d'appeler un setter dans le corps ou de muter un objet d'état.
\reactant{} la traite donc comme une \emph{hypothèse} dont chaque violation a sa
propre règle : \regle{setter-in-render} pour un setter appelé pendant le rendu,
\regle{state-mutation} pour une mutation en place. Ce motif --- une hypothèse du
modèle abstrait, doublée d'une règle qui signale sa violation --- reviendra
plusieurs fois ; le projet l'appelle \emph{soundness en couches}
(\cref{sec:react-modele-identite}).

\subsection{Ce que l'analyseur voit de \code{Counter}}

Le front-end de \reactant{} (\cref{chap:frontend-detection}) détecte
\code{Counter} comme un composant, et le lowering (\cref{chap:lowering-cfg}) en
produit deux choses : un graphe de flot de contrôle du corps, le
\code{render_cfg}, et une liste d'entrées de hooks. Pour \code{Counter}, cette
liste contient un \code{HookEntry::State} (le \code{useState}, label 0) et un
\code{HookEntry::Handler} (le \code{onClick}, label 1) : le handler n'est pas un
hook React, mais l'IR l'extrait comme un \emph{point d'entrée} séparé
(\cref{sec:react-modele-concrete}). D'où le « 2 hooks » de la sortie :

\begin{shell}
$ reactant check --no-color --info --show-clean --project plain counter.tsx
\end{shell}
\begin{sortie}
  Counter  (2 hooks)  counter.tsx  ✓
    verified  conditional-hook  all hooks run unconditionally, in a stable order
    verified  lazy-init  no useState/useRef initializer re-runs work on every render
    verified  setter-in-render  no setter is called during render
    verified  state-mutation  no state or prop object is mutated in place

✓  1 file(s) no issues found.
\end{sortie}

Aucun diagnostic : un setter dans un handler est l'usage normal de React. Ce
silence n'est pas une évidence pour un analyseur. Le handler peut s'exécuter un
nombre quelconque de fois, et rien ne borne la valeur de \code{count} ; il faut
donc que le moteur sache qu'une croissance causée par des clics n'est pas une
boucle infinie. C'est
exactement la distinction « déclenché par le cycle » / « déclenché de
l'extérieur » de la \cref{fig:react-modele-chronologie}, étudiée à la
\cref{sec:react-modele-concrete}.

\begin{piege}
Le corps d'un composant n'est pas « le programme ». Un lecteur qui aborde
l'analyse comme celle d'une fonction JavaScript ordinaire oublie que React
l'appelle en boucle, garde l'état \emph{hors} de la fonction, et exécute d'autres
morceaux du code (effets, handlers) à d'autres moments. Toute la difficulté de
l'analyse --- et tout ce chapitre --- tient dans ce décalage.
\end{piege}

\section{L'état : \code{useState}, la file d'updaters et le batching}
\label{sec:react-modele-usestate}

\subsection{Un setter ne modifie pas la variable}

Le premier malentendu sur React porte sur ce que fait un setter. Le programme
suivant l'expose.

\begin{lstlisting}[language=TypeScript,caption={\code{Batch} : trois appels au setter dans un même handler (\file{/tmp/ch02/batching.tsx})},label={lst:react-modele-batch}]
import { useState } from "react";

export function Batch() {
  const [n, setN] = useState(0);
  function onClick() {
    setN(n + 1);
    setN(n + 1);
    setN((m) => m + 1);
  }
  return <button onClick={onClick}>{n}</button>;
}
\end{lstlisting}

Un lecteur qui lit \code{setN} comme une affectation s'attend à ce qu'un clic
fasse passer \code{n} de \code{0} à \code{3}. Il n'en est rien : \code{n} est
une constante du rendu courant, elle vaut \code{0} pendant tout le handler, et
chaque appel au setter \emph{ajoute une mise à jour à une file}. La file du
premier clic contient trois entrées, que React rejoue au rendu suivant, dans
l'ordre :

\begin{center}\small
\begin{tabular}{@{}lll@{}}
  \toprule
  mise à jour en file & valeur reçue & valeur produite \\
  \midrule
  \og remplacer par \code{0 + 1}\fg{} & \code{0} (ignorée) & \code{1} \\
  \og remplacer par \code{0 + 1}\fg{} & \code{1} (ignorée) & \code{1} \\
  \code{m => m + 1} & \code{1} & \code{2} \\
  \bottomrule
\end{tabular}
\end{center}

Le bouton affiche donc \code{2}. Les deux premières entrées sont des
\emph{remplacements} calculés avec la valeur figée \code{n = 0} ; seule la
troisième, un \terme{updater} (une fonction de l'état précédent vers le
suivant), lit la valeur produite par les mises à jour qui la précèdent.

\begin{react}[set, file d'updaters, batching]
La documentation de \code{useState} \cite{ReactDocs} est explicite :
\og The set function only updates the state variable for the \emph{next}
render. If you read the state variable after calling the set function, you
will still get the old value that was on the screen before your call.\fg{}
Les mises à jour sont regroupées (\emph{batching}) : \og React batches state
updates. It updates the screen after all the event handlers have run and have
called their set functions.\fg{} Pour les updaters, \og React queues this
function to be processed after all the other code in the event handler has
run. During the next render, React goes through the queue and gives you the
final updated state.\fg{} Enfin, le regroupement ne franchit pas les
événements : \og React does not batch across multiple intentional events like
clicks\fg{}. L'exemple canonique de la page \og Queueing a Series of State
Updates\fg{} : \code{setNumber(number + 5)} suivi de
\code{setNumber(n => n + 1)} donne \code{6}.
\end{react}

\begin{piege}
\og \code{set} modifie l'état\fg{} est faux ; \og \code{set} planifie le
prochain rendu\fg{} est juste. Deux conséquences qui trompent régulièrement
les contributeurs de l'analyseur. D'abord, lire l'état juste après l'avoir
écrit, dans le même corps, donne l'\emph{ancienne} valeur : un modèle
abstrait qui ferait \og voir\fg{} la nouvelle valeur à la ligne suivante
serait faux (nous verrons que \reactant{} voit l'ancienne \emph{et} la
nouvelle, ce qui est permis). Ensuite, \code{setN(n + 1)} et
\code{setN((m) => m + 1)} ne sont pas interchangeables : le premier capture
\code{n} au moment où la fermeture a été créée, le second lit la file. C'est
tout le ressort de la \emph{stale closure} du \code{Timer} de l'introduction
(\cref{lst:introduction-timer}).
\end{piege}

\subsection{Le bail-out \code{Object.is}}

Un setter qui écrit la valeur déjà présente ne relance pas le rendu.

\begin{react}[bail-out]
\og If the new value you provide is identical to the current state, as
determined by an \code{Object.is} comparison, React will skip re-rendering the
component and its children.\fg{} \cite{ReactDocs} (page \code{useState}).
\end{react}

Ce détail décide de la terminaison de nombreux programmes. Dans le composant
suivant, l'effet n'a pas de deps et tourne donc après chaque commit ; mais il
écrit \code{0} dans un état qui vaut déjà \code{0}, React compare, ne rend pas,
et il n'y a pas de second commit.

\begin{lstlisting}[language=TypeScript,caption={\code{SameValue} : un effet qui réécrit la valeur courante (\file{/tmp/ch02/samevalue.tsx}, imports omis)},label={lst:react-modele-samevalue}]
export function SameValue() {
  const [n, setN] = useState(0);
  useEffect(() => {
    setN(0);
  });
  return <div>{n}</div>;
}
\end{lstlisting}

\begin{sortie}
  SameValue  (2 hooks)  samevalue.tsx
    warn   redundant-set-state  [hook:0]  (line 5:2)  state `n` is set to a stable value it already holds, so the update is redundant
\end{sortie}

L'analyseur ne signale pas de boucle (la valeur abstraite de \code{n} reste le
point $[0,0]$ et ne croît pas), et la règle \regle{redundant-set-state} relève
l'écriture inutile.

\begin{remarque}
Conventions pour les sorties de ce chapitre. Toutes ont été obtenues au commit
\code{e67b10a} par \code{reactant check --no-color --info --show-clean
--project plain <fichier>} (le binaire de \file{target/debug}). Elles sont
élaguées, sans que ce soit redit, des lignes \code{verified} qui ne portent pas
l'information, des lignes de bilan final et des lignes
\og (\emph{k} trace step(s), rerun with --trace)\fg{} ; quand une sortie
montre des lignes en flèche (\code{→}), celles-ci viennent de la même
commande relancée avec \code{--trace}. Les lignes \code{[verbose]} viennent de
\code{--verbose}.
\end{remarque}

\subsection{L'état initial et l'initialiseur paresseux}

L'argument de \code{useState} n'est lu qu'une fois.

\begin{react}[état initial]
\og React saves the initial state once and ignores it on the next
renders.\fg{} Une fonction passée en argument est un \emph{initialiseur} :
\og If you pass a function as initialState, it will be treated as an
initializer function. It should be pure, should take no arguments, and should
return a value of any type.\fg{} La documentation oppose
\code{useState(createInitialTodos)} (la fonction est passée, appelée une seule
fois) à \code{useState(createInitialTodos())} (l'appel est évalué à chaque
rendu, son résultat jeté après le premier) \cite{ReactDocs}.
\end{react}

\begin{lstlisting}[language=TypeScript,caption={\code{Todos} : initialiseur appelé à chaque rendu, et sa forme paresseuse (\file{/tmp/ch02/lazy.tsx}, imports omis)},label={lst:react-modele-lazy}]
function createTodos() {
  return [{ id: 1, done: false }];
}

export function Todos() {
  const [todos] = useState(createTodos());
  const [lazy] = useState(() => createTodos());
  return <ul>{todos.length + lazy.length}</ul>;
}
\end{lstlisting}

\begin{sortie}
  Todos  (2 hooks)  lazy.tsx
    warn   lazy-init  [hook:0]  (line 8:8)  this useState is initialised by a direct function call. The call runs on every render but the result is only used on mount; wrap as `useState(() => ...)` to defer it
\end{sortie}

\subsection{Ce que fait le code}

\paragraph{Le lowering.} Chaque appel de hook reconnu comme venant de React
devient une entrée \code{HookEntry} (la liste complète est au
\cref{sec:react-modele-regles-hooks}). Pour \code{useState}, l'entrée ne garde
que l'expression initiale :

\rustsnippet{src/lowering/hook_extractor.rs}{667}{681}{un \code{useState} devient \code{HookEntry::State}, un \code{useEffect} un \code{HookEntry::Effect}}

La valeur lue de l'état, \code{n}, devient dans le CFG du rendu l'expression
\code{Expr::StateVal(label)} et le setter \code{Expr::StateSetter(label)} ;
le label est la position du hook (\cref{sec:react-modele-regles-hooks}).

\paragraph{L'amorçage.} Avant la boucle de point fixe, le moteur évalue une fois
chaque initialiseur et en fait la valeur de départ du slot. Un initialiseur
paresseux (une fonction sans paramètre) est \emph{exécuté}, et c'est sa valeur
de retour qui amorce le slot, conformément à la documentation :

\rustsnippet{src/engine/fixpoint.rs}{336}{350}{amorçage d'un slot, initialiseur paresseux compris}

Cet amorçage a lieu \emph{hors} de la boucle (\src{src/engine/fixpoint.rs}{314}{353}) :
l'initialiseur ne ré-influence jamais le store, exactement comme React
\og ignores it on the next renders\fg{}. La règle \regle{lazy-init} exploite
l'autre moitié du fait React : l'expression non paresseuse, elle, est
ré-évaluée à chaque rendu.

\paragraph{Le store d'état.} L'état abstrait d'un composant est un
\code{StateStore}, une table de label vers valeur abstraite :

\rustsnippet{src/domains/stores/state_store.rs}{24}{33}{lecture et mise à jour du store d'état}

Tout tient dans la seconde fonction. Écrire dans un slot ne \emph{remplace}
pas sa valeur : on en prend le join avec la valeur écrite. C'est une
\terme{mise à jour faible} (\emph{weak update}), et le store ne fait que
croître selon l'ordre du treillis. Un slot jamais écrit ni amorcé vaut
$\Bot$.

\paragraph{L'exécution d'un setter.} Quand l'interprète abstrait rencontre un
appel \code{setX(arg)}, il exécute \code{exec_setter_call}, dont le
commentaire de tête (\src{src/domains/interp/interpreter.rs}{335}{340}) dit
l'intention : \og weak-update the corresponding state label\fg{}, en traitant
les updaters fonctionnels par l'exécution de leur corps avec le paramètre lié
à la valeur courante de l'état.

\rustsnippet[label={lst:react-modele-exec-setter}]{src/domains/interp/interpreter.rs}{341}{366}{\code{exec_setter_call} : un setter est un join, un updater est exécuté sur le store}

Lisons-le branche par branche.
\begin{itemize}
  \item La garde (l.~348--350) reconnaît un appel dont le callee est une
    variable liée à un setter de ce composant (\code{env.setter_label}).
  \item Si l'argument est une fonction littérale (l.~353--361), c'est un
    updater : son paramètre est lié à \code{ctx.state.get(label)}, la valeur
    \emph{courante du store}, et son corps est exécuté ; la valeur écrite est
    ce que l'updater renvoie.
  \item Sinon (l.~362), l'argument est évalué dans l'environnement courant ;
    sans argument (l.~363), on écrit $\Top$.
  \item Dans tous les cas, l.~365, la valeur est jointe au slot.
\end{itemize}
La suite de la fonction (\src{src/domains/interp/interpreter.rs}{368}{392})
traite l'appel d'un setter \emph{d'un autre composant}, en analyse
inter-composants : l'argument est évalué et joint dans un
\code{SharedStateStore} indexé par (composant, label). Nous y reviendrons au
\cref{sec:react-modele-concrete}, car cette branche est aussi le siège d'un
défaut vérifié.

\begin{exemple}{\code{Batch} dans le domaine abstrait}{usestate-batch}
Rejouons à la main la passe du handler de \code{Batch} (\cref{lst:react-modele-batch}) lors de
la première itération, en ne suivant que l'intervalle numérique. Le store vaut
$\{n \mapsto [0,0]\}$ et, dans l'environnement du handler, \code{n} vaut
$[0,0]$.
\begin{enumerate}
  \item \code{setN(n + 1)} : on écrit $[1,1]$ ; le store devient
    $[0,0] \join [1,1] = [0,1]$.
  \item \code{setN(n + 1)} : \code{n} est lu dans l'environnement, qui n'a pas
    changé ($[0,0]$) ; on écrit encore $[1,1]$ ; le store reste $[0,1]$.
  \item \code{setN((m) => m + 1)} : \code{m} est lié au store, $[0,1]$ ;
    l'updater renvoie $[1,2]$ ; le store devient $[0,2]$.
\end{enumerate}
React produit \code{2} ; l'analyseur a calculé $[0,2]$, qui contient \code{2}
mais aussi \code{0} et \code{1}. Il ne représente ni la file, ni son ordre, ni
le moment où elle est rejouée : il joint tout, tout de suite. (Ce calcul est
celui de la première branche d'\code{exec_setter_call}, seule active en analyse
intra ; sur le chemin programme de la CLI, la seconde branche s'y ajoute pour
l'updater, avec l'effet décrit au \cref{sec:react-modele-concrete}.)
\end{exemple}

Ce calcul est imprécis, mais il est \emph{sûr} : tout état concret
atteignable après n'importe quelle suite de clics et de rendus est obtenu en
appliquant des mises à jour à des valeurs que le store contient déjà, et le
store, fermé par join, contient leurs résultats. L'argument est développé au
\cref{sec:react-modele-concrete} ; il repose sur deux choix que l'on retrouvera
partout : la mise à jour est un join, et l'updater lit un sur-ensemble de la
valeur que React lui passerait.

\begin{remarque}
Lire l'état après l'avoir écrit dans le même corps donne, dans l'analyseur,
l'ancienne valeur \emph{jointe} à la nouvelle, alors que React donne
l'ancienne. Le join contient la réponse de React : le résultat est imprécis
mais correct. C'est le sens du premier piège de la section.
\end{remarque}

Le bail-out \code{Object.is}, lui, n'est pas modélisé dans le point fixe : le
moteur suppose que tout set peut relancer un rendu. C'est encore une
sur-approximation (on considère plus de rendus que React n'en fait), et le fait
React n'est exploité que par des règles qui en ont besoin, comme
\regle{redundant-set-state} ci-dessus ou \regle{state-mutation}, dont la
documentation rappelle que muter un objet en place laisse \code{Object.is}
conclure à l'égalité et saute le rendu
(\src{src/rules/impls/state_mutation.rs}{21}{26}).

\paragraph{\code{useReducer}.} Le lowering range \code{useReducer} dans la même
case que \code{useState} :

\rustsnippet{src/lowering/hook_extractor.rs}{715}{719}{\code{useReducer} : le réducteur est jeté}

Le réducteur n'est jamais exécuté. \code{dispatch(action)} est donc lu comme un
setter qui écrit \emph{l'action} dans le slot. Le \cref{sec:react-modele-concrete}
montre la conséquence observée : une \Error{} sur un programme qui ne boucle
pas.

\section{Les effets : après le commit, gardés par les deps}
\label{sec:react-modele-effets}

\subsection{Quand un effet s'exécute}

Un effet est la seule partie du composant que React exécute \emph{après} le
commit. C'est aussi la seule qui puisse, à elle seule, relancer un tour sans
intervention extérieure : c'est donc elle qui porte les boucles infinies.

\begin{react}[\code{useEffect}]
Signature : \code{useEffect(setup, dependencies?)}. \og React will compare each
dependency with its previous value using the \code{Object.is} comparison. If
you omit this argument, your Effect will re-run after every commit of the
component.\fg{} Avec \code{[]}, l'effet ne tourne qu'après le commit initial.
\og After every commit with changed dependencies, React will first run the
cleanup function (if you provided it) with the old values, and then run your
setup function with the new values\fg{} ; au démontage, le cleanup tourne une
dernière fois. Le moment exact : \og If your Effect wasn't caused by an
interaction (like a click), React will generally let the browser paint the
updated screen first before running your Effect.\fg{} Enfin, \og Effects only
run on the client\fg{} \cite{ReactDocs}.
\end{react}

Les trois formes du tableau de deps donnent trois comportements, que le
programme suivant met côte à côte avec le même corps d'effet.

\begin{lstlisting}[language=TypeScript,caption={Trois formes de deps pour le même effet (\file{/tmp/ch02/deps.tsx}, imports omis)},label={lst:react-modele-deps}]
export function EveryCommit() {
  const [n, setN] = useState(0);
  useEffect(() => {
    setN(n + 1);
  });
  return <div>{n}</div>;
}

export function MountOnly() {
  const [n, setN] = useState(0);
  useEffect(() => {
    setN(n + 1);
  }, []);
  return <div>{n}</div>;
}

export function OnChange() {
  const [n, setN] = useState(0);
  useEffect(() => {
    setN(n + 1);
  }, [n]);
  return <div>{n}</div>;
}
\end{lstlisting}

À l'exécution :
\begin{itemize}
  \item \code{EveryCommit} boucle : chaque commit exécute l'effet, qui écrit
    une valeur nouvelle, donc un nouveau rendu, donc un nouveau commit ;
  \item \code{MountOnly} rend deux fois (\code{0} puis \code{1}) et s'arrête :
    l'effet ne tourne qu'une fois ;
  \item \code{OnChange} boucle aussi : \code{n} change à chaque tour, et c'est
    justement sa dep.
\end{itemize}

\begin{sortie}
  EveryCommit  (2 hooks)  deps.tsx
    warn   infinite-loop  [hook:0]  (line 5:2)  this effect keeps pushing state `n` to new values on every run. Potential infinite render loop
    info   widening-info  (line 5:2)  state `n` kept changing during analysis and was approximated to converge, so findings that depend on it may be imprecise
  MountOnly  (2 hooks)  deps.tsx
    warn   missing-deps  [hook:1]  var:n  (line 13:2)  `n` is used in this effect but not in its deps array, and it is recreated on every render
    info   widening-info  (line 13:2)  state `n` kept changing during analysis and was approximated to converge, so findings that depend on it may be imprecise
    verified  infinite-loop  no effect diverges into an infinite render loop
  OnChange  (2 hooks)  deps.tsx
    warn   infinite-loop  [hook:0]  (line 21:2)  this effect keeps pushing state `n` (its deps do not provably gate it, so the effect can re-run every render) to new values on every run. Potential infinite render loop
    info   widening-info  (line 21:2)  state `n` kept changing during analysis and was approximated to converge, so findings that depend on it may be imprecise
\end{sortie}

Les deux boucles sont signalées en \Warning{}, jamais en \Error{} : la
divergence d'une valeur numérique ne donne pas, à ce jour, la preuve complète
qu'exige le niveau \Error{} (\issue{144}). \code{MountOnly} n'a pas de boucle,
mais un \Info{} \regle{widening-info} y apparaît : le piège de la fin de section
explique pourquoi.

\subsection{\code{SelfCounter} et \code{Inf} : le cycle Check--Effect}

L'article React-tRace \cite{LeeAhnYi2025} illustre l'effet par deux composants
presque identiques. Leur transcription TSX :

\begin{exemple}{\code{Inf} et \code{SelfCounter}}{canon-selfcounter}
\begin{lstlisting}[language=TypeScript,caption={\code{Inf} et \code{SelfCounter} (\file{/tmp/ch02/selfcounter.tsx})},label={lst:react-modele-selfcounter}]
import { useEffect, useState } from "react";

export function Inf() {
  const [s, setS] = useState(0);
  useEffect(() => {
    setS(s + 1);
  });
  return <div>{s}</div>;
}

export function SelfCounter() {
  const [s, setS] = useState(0);
  useEffect(() => {
    if (s < 3) setS(s + 1);
  }, [s]);
  return <div>{s}</div>;
}
\end{lstlisting}
\end{exemple}

\code{SelfCounter} rend quatre fois de suite sans aucune interaction
(\code{s} vaut \code{0}, \code{1}, \code{2}, \code{3}) puis s'arrête, comme le
montre la \cref{fig:react-modele-selfcounter} ; \code{Inf}, privé de la garde,
ne s'arrête jamais. L'article résume le mécanisme : \og Infinite render loop
occurs because setting state within an Effect triggers the runtime to Check if
the state has changed [\dots]. If the state whose setter function is called is
actually modified, the component re-renders and decides to run the Effect
again. Essentially, this Check-Effect decision cycle creates a render
loop.\fg{} Les mots \emph{Check} et \emph{Effect} sont ici des termes
techniques, les \emph{décisions} de la sémantique formelle
(\cref{sec:react-modele-trace}).

\ifhastikz
\begin{figure}[htbp]\centering
\begin{tikzpicture}[x=2.55cm,y=1cm,
    tour/.style={bloc,minimum width=20mm,font=\scriptsize}]
  \foreach \i/\s/\act in {0/0/{set 1},1/1/{set 2},2/2/{set 3},3/3/{aucun set}} {
    \node[tour] (r\i) at (\i,0) {rendu\\\code{s = \s}};
    \node[tour,below=5mm of r\i] (e\i) {effet\\\act};
    \draw[fleche] (r\i) -- node[etiquette,right]{commit} (e\i);
  }
  \foreach \i/\j in {0/1,1/2,2/3} {
    \draw[flechemust] (e\i.east) -- (r\j.west);
  }
  \node[bloc,right=6mm of e3,font=\scriptsize] (loop) {boucle\\d'événements};
  \draw[fleche] (e3) -- (loop);
  \node[etiquette,above=1mm of r0] {montage};
  \node[etiquette,above=1mm of r1] {\code{Object.is(1, 0)} faux};
  \node[etiquette,above=1mm of r2] {\code{Object.is(2, 1)} faux};
  \node[etiquette,above=1mm of r3] {\code{Object.is(3, 2)} faux};
\end{tikzpicture}
\caption{Déroulement concret de \code{SelfCounter}. Chaque rendu est suivi d'un
commit puis de l'effet, car la dep \code{s} a changé. Le quatrième effet tourne
encore (la dep a changé de \code{2} à \code{3}) mais n'écrit rien : aucun tour
n'est relancé et React attend un événement. Pour \code{Inf}, la chaîne ne
s'arrête pas.}
\label{fig:react-modele-selfcounter}
\end{figure}
\fi

\begin{sortie}
  Inf  (2 hooks)  selfcounter.tsx
    warn   infinite-loop  [hook:0]  (line 5:2)  this effect keeps pushing state `s` to new values on every run. Potential infinite render loop
       → state `s` is written here [hook:1] (line 5:2)
       → the abstract value of state `s` kept growing and was widened at iteration 3
    info   widening-info  (line 5:2)  state `s` kept changing during analysis and was approximated to converge, so findings that depend on it may be imprecise
  SelfCounter  (2 hooks)  selfcounter.tsx
    info   widening-info  (line 13:2)  state `s` kept changing during analysis and was approximated to converge, so findings that depend on it may be imprecise
    verified  infinite-loop  no effect diverges into an infinite render loop
\end{sortie}

La sortie (obtenue avec \code{--trace} pour les lignes en flèche) raconte le
calcul abstrait. Pour \code{Inf}, le slot \code{s} vaut $[0,0]$, puis
$[0,1]$, puis $[0,2]$ ; au troisième tour de la boucle abstraite, le moteur
applique un \emph{widening} et le slot passe à $[0,+\infty)$ ; ce slot entre
dans la trace d'élargissement (\code{widen_trace}), et la règle
\regle{infinite-loop} en conclut que l'effet pousse l'état vers des valeurs
toujours nouvelles. Pour \code{SelfCounter}, le slot est aussi élargi (d'où
l'\Info), mais l'écriture elle-même, recalculée après convergence, reste
bornée par la garde \code{s < 3} : pas de boucle. Le mécanisme exact
(widening à seuils, écritures « pures » des effets) est l'objet du
\cref{chap:point-fixe-composant}.

\subsection{\code{Flicker} : un rendu de trop}

Un effet de montage qui écrit une constante ne boucle pas, mais coûte un rendu :
la valeur initiale s'affiche une fraction de seconde. L'article appelle ce
composant \code{Flicker} (\og A swift user might even see the transient state
0\fg{}).

\begin{exemple}{\code{Flicker}}{canon-flicker}
\begin{lstlisting}[language=TypeScript,caption={\code{Flicker} (\file{/tmp/ch02/flicker.tsx}, imports omis)},label={lst:react-modele-flicker}]
export function Flicker() {
  const [s, setS] = useState(0);
  useEffect(() => {
    setS(42);
  }, []);
  return <div>{s}</div>;
}
\end{lstlisting}
\begin{sortie}
  Flicker  (2 hooks)  flicker.tsx
    warn   unnecessary-rerender  [hook:0]  (line 5:2)  mount-only effect sets state `s` to a constant different from its initial value, which costs one extra rerender on mount; consider initialising directly with the target value
\end{sortie}
\end{exemple}

\subsection{\code{Toggle} : une boucle que la valeur abstraite ne voit pas}

Le dernier exemple canonique de cette section est un avertissement pour le
lecteur, et une limite de l'outil au commit étudié.

\begin{exemple}{\code{Toggle}}{canon-toggle}
\begin{lstlisting}[language=TypeScript,caption={\code{Toggle} (\file{/tmp/ch02/toggle.tsx})},label={lst:react-modele-toggle}]
import { useEffect, useState } from "react";

export function Toggle() {
  const [on, setOn] = useState(false);
  useEffect(() => {
    setOn(!on);
  }, [on]);
  return <div>{on ? 1 : 0}</div>;
}
\end{lstlisting}
\end{exemple}

En React, ce composant boucle : \code{on} passe de \code{false} à
\code{true}, la dep a changé, l'effet repasse, \code{on} revient à
\code{false}, et ainsi de suite, un rendu par tour. Pourtant l'ensemble des
valeurs atteintes est fini : $\{\code{false}, \code{true}\}$. L'analyseur,
qui calcule précisément cet ensemble, converge sans widening :

\begin{sortie}
  [verbose] Toggle: 1 iteration(s), widened: []
✓  1 file(s) no issues found.
\end{sortie}

\begin{piege}
Un point fixe abstrait atteint ne prouve pas que la boucle de React s'arrête.
La sémantique que le moteur sur-approxime est l'ensemble des \emph{états}
atteignables, pas la longueur des suites de rendus. Une boucle qui oscille
entre des valeurs déjà vues (\code{Toggle}) a un ensemble d'états fini ; sa
détection ne peut venir que d'un argument de \emph{structure} (la dep est le
slot écrit, et la valeur écrite diffère toujours de la valeur lue), du même
genre que celui du churn d'objets (\cref{sec:react-modele-identite}). Au commit
\code{e67b10a}, ce raisonnement existe pour les références fraîches, pas pour
une négation booléenne : \code{Toggle} est un faux négatif observé.
\end{piege}

\subsection{Cleanup et effets de mise en page}

Un effet peut renvoyer une fonction de nettoyage. React l'appelle avant chaque
ré-exécution et au démontage. Un effet qui s'abonne sans se désabonner
empile les abonnements :

\begin{lstlisting}[language=TypeScript,caption={\code{Width} : un abonnement sans cleanup (\file{/tmp/ch02/listener.tsx}, imports omis)},label={lst:react-modele-listener}]
export function Width() {
  const [w, setW] = useState(0);
  const box = useRef(null);
  useEffect(() => {
    window.addEventListener("resize", () => setW(window.innerWidth));
  }, [box]);
  return <div ref={box}>{w}</div>;
}
\end{lstlisting}

\begin{sortie}
  Width  (4 hooks)  listener.tsx
    warn   missing-cleanup  [hook:2]  (line 7:4)  this effect calls `window.addEventListener` but returns no cleanup. The registration is repeated every time the effect re-runs (and on every mount, twice under StrictMode) and nothing ever undoes it; return a function that tears it down
    verified  infinite-loop  no effect diverges into an infinite render loop
\end{sortie}

Deux détails de la sortie comptent pour ce chapitre. Les « 4 hooks » sont le
\code{useState}, le \code{useRef}, l'effet et le callback d'\code{addEventListener},
extrait comme un handler (\code{HookEntry::Handler}) par le lowering
(\code{extract_subscriptions}, \src{src/lowering/hook_extractor.rs}{16}{26}) : il ne tourne que sur un
événement extérieur. Et \regle{infinite-loop} est vérifiée : l'écriture
\code{setW} n'a lieu que sur un redimensionnement, pas dans le cycle
rendu--effet (\cref{sec:react-modele-concrete}).

\code{useLayoutEffect} et \code{useInsertionEffect} s'exécutent plus tôt que
\code{useEffect} (avant la peinture). Pour l'analyseur, la différence de moment
ne change rien à ce qui compte --- l'effet ne s'exécute jamais pendant le rendu,
et il s'exécute après le commit --- et le lowering les confond :

\rustsnippet{src/lowering/hook_extractor.rs}{720}{729}{\code{useLayoutEffect} et \code{useInsertionEffect} sont des \code{HookEntry::Effect}}

\subsection{Ce que fait le code}

\paragraph{Tous les effets, à chaque itération.} La boucle de point fixe
(\cref{chap:point-fixe-composant}) exécute, à chaque itération, le rendu puis
chaque effet, sur le même instantané du store et dans l'environnement de sortie
du rendu :

\rustsnippet[label={lst:react-modele-effect-pass}]{src/engine/fixpoint.rs}{415}{449}{les passes d'effets de la boucle de point fixe}

Le motif \code{HookEntry::Effect { label, body_cfg, .. }} de la l.~419
ignore le champ \code{deps} : \emph{tous} les effets sont exécutés à chaque
itération, quels que soient leurs deps. C'est la décision de modélisation la
plus importante de la section.

\begin{decision}[tous les effets à chaque tour abstrait]
React exécute un effet \emph{si et seulement si} l'une de ses deps a changé
selon \code{Object.is} (ou s'il n'a pas de deps, ou au montage). Le point fixe
l'exécute \emph{toujours}. L'ensemble des écritures possibles en est un
sur-ensemble : c'est sound. La précision perdue est récupérée dans les règles,
qui relisent les deps : \regle{infinite-loop} écarte les effets de montage et
ceux dont toutes les deps sont prouvées stables
(\src{src/rules/impls/infinite_loop.rs}{96}{113}), et la relation
\relation{effect_triggers} dit quand un effet \emph{doit} repartir. L'ADR-004
(\adr{4}) décrivait une passe « for each Effect whose decision = Effect » ; le
code n'a pas de notion de décision, et c'est plus simple et plus sûr.
\end{decision}

Chaque effet démarre du même store et du même environnement (celui du rendu de
cette itération) : l'ordre dans lequel React exécute les effets (enfants
d'abord, puis ordre syntaxique) est abstrait par le join, qui est commutatif ;
l'influence d'un effet sur un autre est vue à l'itération suivante. Et un effet
voit les valeurs \emph{du rendu qui l'a créé} : c'est une fermeture, comme dans
React.

\paragraph{Quand un effet doit repartir.} Pour transformer « l'effet
\emph{peut} repartir » en « l'effet \emph{doit} repartir », il faut savoir
qu'une dep change à coup sûr. Le moteur calcule pour cela une relation, une
ligne par couple (dep, slot) :

\rustsnippet{src/engine/triggers.rs}{33}{44}{une ligne de la relation \relation{effect_triggers}}

Le bit \code{exact} est le cœur : la dep \emph{est} la valeur du slot (ou un
alias), donc si l'on stocke dans ce slot une valeur fraîche, \code{Object.is}
échoue sur cette dep et l'effet repart. Sans ce bit, la dep est seulement
\emph{versionnée} par le slot (un champ, un mémo, une prop calculée par le
parent) et l'effet \emph{peut} repartir. Le commentaire de module précise
l'intention : \og This answers identity, not dependence\fg{}
(\src{src/engine/triggers.rs}{1}{17}). La relation est étudiée au
\cref{chap:relations} ; retenons qu'elle est la traduction abstraite de la
phrase « React compare chaque dep avec \code{Object.is} ».

\begin{piege}
\code{widen_trace} ne signifie pas « boucle React ». Dans \code{MountOnly}, le
slot est élargi (l'\Info{} \regle{widening-info} le dit) parce que le point fixe
exécute l'effet à chaque itération, deps \code{[]} ou non. La trace
d'élargissement dit « la valeur abstraite croît sous l'hypothèse que tout effet
peut toujours repartir » ; c'est la règle, en relisant les deps, qui décide
s'il y a boucle. Lire \code{widen_trace} seul produirait un faux positif sur
chaque effet de montage.
\end{piege}

\section{L'identité référentielle}
\label{sec:react-modele-identite}

\subsection{Même contenu, autre objet}

React compare les deps, l'état et la valeur d'un contexte avec
\code{Object.is}. Pour un nombre, une chaîne ou un booléen, c'est une
comparaison de valeur ; pour un objet, un tableau ou une fonction, c'est une
comparaison d'\emph{identité} : deux objets de même contenu sont différents
s'ils ont été alloués séparément.

\begin{react}[identité et mémoïsation]
\og In JavaScript, a \code{function () {}} or \code{() => {}} always creates
a \emph{different} function, similar to how the \code{{}} object literal always
creates a new object\fg{} (page \code{useCallback}). La documentation de
\code{useEffect} en tire la conséquence : un objet construit dans le corps du
composant et placé dans les deps fait repartir l'effet à chaque rendu ; le
remède est de le construire \emph{dans} l'effet, ou de le mémoïser.
\code{useMemo} et \code{useCallback} renvoient la valeur en cache tant que
leurs deps sont égales au sens de \code{Object.is}, et la recalculent sinon ;
sans tableau, à chaque rendu. \code{useRef} renvoie le même objet à tous les
rendus, et modifier \code{ref.current} ne déclenche pas de rendu
\cite{ReactDocs}.
\end{react}

\begin{lstlisting}[language=TypeScript,caption={\code{Search} : une dep fraîche à chaque rendu, et sa version mémoïsée (\file{/tmp/ch02/search.tsx}, imports omis)},label={lst:react-modele-search}]
export function Search({ query }: { query: string }) {
  const [hits, setHits] = useState(0);
  const options = { query, limit: 10 };
  const stable = useMemo(() => ({ query, limit: 10 }), [query]);
  useEffect(() => {
    console.log(options.limit);
  }, [options]);
  useEffect(() => {
    console.log(stable.limit, hits);
  }, [stable, hits]);
  return <button onClick={() => setHits(hits + 1)}>{hits}</button>;
}
\end{lstlisting}

\begin{sortie}
  Search  (5 hooks)  search.tsx
    warn   always-unstable-deps  [hook:2]  (line 7:2)  this effect depends on `options`, a new reference every render, so `Object.is` always differs and the effect re-runs on every render regardless of the other deps
\end{sortie}

Le premier effet repart à chaque rendu : \code{options} est un littéral objet,
donc un objet neuf à chaque appel du corps. Le second ne repart que si
\code{query} ou \code{hits} change : \code{stable} est la même référence tant
que \code{query} ne bouge pas. Les « 5 hooks » sont l'état, le mémo, les deux
effets et le handler du bouton.

\subsection{\code{ObjChurn} : une boucle par l'identité}

L'identité peut aussi fabriquer une boucle infinie sans qu'aucune valeur ne
« croisse ». C'est l'exemple qui a motivé l'ADR-017 (\adr{17}).

\begin{exemple}{\code{ObjChurn} et \code{ObjStateDep}}{canon-objchurn}
\begin{lstlisting}[language=TypeScript,caption={\code{ObjChurn} et \code{ObjStateDep} (\file{/tmp/ch02/objchurn.tsx})},label={lst:react-modele-objchurn}]
import { useEffect, useState } from "react";

export function ObjChurn() {
  const [obj, setObj] = useState({ a: 1 });
  useEffect(() => {
    setObj({ ...obj, b: 2 });
  }, [obj]);
  return <div>{obj.a}</div>;
}

export function ObjStateDep() {
  const [ctx] = useState({ locale: "en" });
  useEffect(() => {
    document.title = ctx.locale;
  }, [ctx]);
  return <div>{ctx.locale}</div>;
}
\end{lstlisting}
\end{exemple}

Dans \code{ObjChurn}, l'effet écrit un objet neuf ; \code{Object.is} échoue
sur la dep \code{obj} ; l'effet repart, et ainsi de suite : c'est une boucle
certaine, alors que le \emph{contenu} de l'objet est le même dès le deuxième
tour. Dans \code{ObjStateDep}, l'état est aussi un objet littéral, mais rien ne
l'écrit : sa référence est la même à chaque rendu, et l'effet ne tourne qu'au
montage.

\begin{sortie}
  ObjChurn  (2 hooks)  objchurn.tsx
    error  infinite-loop  [hook:0]  (line 5:2)  this effect recreates object state `obj` it depends on. Every run stores a fresh reference (`Object.is` always fails) and re-triggers itself: infinite render loop
       → a fresh value is written to state `obj` here [hook:1] (line 6:4)
  ObjStateDep  (2 hooks)  objchurn.tsx  ✓
    verified  always-unstable-deps  no deps array is defeated by an always-fresh reference
    verified  infinite-loop  no effect diverges into an infinite render loop
\end{sortie}

Les deux verdicts sont justes, et chacun demande un fait différent. Pour
\code{ObjChurn}, la divergence de valeur (le mécanisme de \code{Inf}) ne sert à
rien : la valeur abstraite de la référence ne croît pas, \og fraîche\fg{} joint
\og fraîche\fg{} reste \og fraîche\fg{}, et le point fixe converge dès
le premier test (\code{--verbose} : \code{ObjChurn: 0 iteration(s), widened: []}). Pour \code{ObjStateDep}, il faut au contraire savoir que
\emph{lire} un état objet ne donne pas une référence fraîche à chaque rendu,
même si l'état a été initialisé par un littéral.

\subsection{Ce que fait le code}

\paragraph{Les allocations.} L'évaluateur abstrait donne à chaque forme
d'expression sa stabilité. Les littéraux objet, tableau et fonction sont frais
à chaque rendu ; un élément JSX natif, un \code{useRef} sont stables ; un hook
non résolu est $\Top$ :

\rustsnippet{src/domains/transfer/state_value.rs}{142}{156}{l'identité des marqueurs de hooks et des littéraux}

S'y ajoutent \code{new X()} (\srcline{src/domains/transfer/state_value.rs}{181})
et les appels dont le résultat est prouvé neuf, comme \code{map}, \code{filter}
ou \code{JSON.parse} (\code{returns_fresh_reference},
\src{src/domains/transfer/state_value.rs}{238}{259}). La liste exclut à dessein
\code{slice} et \code{concat}, qui renvoient une primitive sur une chaîne, et
les méthodes en place (\code{sort}, \code{reverse}\dots) qui renvoient le
receveur.

\paragraph{Le treillis de stabilité.} La stabilité est une borne sur la
\emph{trace de changement} d'une valeur, l'ensemble des rendus où
\code{Object.is} échoue entre deux rendus consécutifs, qui est la seule chose
que React observe :

\rustsnippet[label={lst:react-modele-stability}]{src/domains/impls/stability.rs}{13}{52}{le treillis \code{Stability}}

La \cref{fig:react-modele-stability} redessine le diagramme du commentaire. Le
treillis est défini et étudié au \cref{chap:statevalue-stabilite}
(\cref{def:stabilite}) ; ce qui nous importe ici est la lecture React
de chaque point. \code{Stable} : la même référence à chaque rendu.
\code{PerRender} : une référence neuve à chaque rendu, \emph{à coup sûr} ---
c'est une borne must, qui permet de tirer. \code{Versioned(S)} : ne change
\emph{qu'aux} événements de set des slots de $S$ --- c'est une borne may, qui
permet de se taire. \code{PerRender} et \code{Versioned} ne sont pas opposés :
ce sont deux bornes différentes, incomparables, dont le join ne garantit rien
(\code{Unknown}).

\ifhastikz
\begin{figure}[htbp]\centering
\begin{tikzpicture}[treillis,node distance=9mm and 14mm]
  \node (top) {\code{Unknown} ($\Top$)};
  \node[below left=of top] (vtop) {\code{VersionedTop}};
  \node[below right=of top] (per) {\code{PerRender}};
  \node[below=of vtop] (vs) {\code{Versioned(S)}};
  \node[below=of vs] (st) {\code{Stable}};
  \node[below right=of st] (bot) {\code{Bottom} ($\Bot$)};
  \draw (top) -- (vtop); \draw (top) -- (per);
  \draw (vtop) -- node[etiquette,left]{$S \subseteq S'$, seuil 4} (vs);
  \draw (vs) -- (st); \draw (st) -- (bot); \draw (per) |- (bot);
  \node[etiquette,left=2mm of vs] {borne may};
  \node[etiquette,right=2mm of per] {borne must};
\end{tikzpicture}
\caption{Le treillis \code{Stability} (ADR-017). À gauche, la chaîne may
« ne change qu'aux sets de $S$ », dont la hauteur est bornée par le seuil
\code{VERSIONED_LABELS_THRESHOLD = 4} ; à droite, le fait must « neuf à chaque
rendu ». Leur join est $\Top$.}
\label{fig:react-modele-stability}
\end{figure}
\fi

\paragraph{La lecture d'un état.} Voici le fait React qui règle
\code{ObjStateDep} : un état ne change \emph{que} lorsque son setter est
appelé. Le moteur le traduit à un seul endroit, la lecture d'un slot :

\rustsnippet[label={lst:react-modele-stateval}]{src/domains/transfer/state_value.rs}{122}{134}{la conversion « côté lecture » de l'ADR-017}

Le store contient le join des valeurs \emph{écrites} ; pour \code{ObjStateDep},
c'est la référence \code{PerRender} du littéral d'initialisation. Mais ce qu'un
rendu \emph{lit} ne peut changer qu'aux événements de set de ce slot : la
lecture remplace donc la composante référence par
\code{Versioned({(composant, label)})}. Une dep \code{Versioned} ne déclenche
pas \regle{always-unstable-deps}, qui exige \code{PerRender}.

\begin{invariant}{Double vue de l'état}{react-modele-double-vue}
Le store d'état est la \emph{vue événement} : pour chaque slot, le join des
valeurs écrites. L'évaluation de \code{Expr::StateVal} est la \emph{vue
inter-rendus} : ce que \code{Object.is} compare d'un rendu à l'autre, dont la
composante référence est toujours \code{Versioned} par le slot lui-même (sauf
si elle est $\Bot$). Les règles qui comparent des valeurs écrites
(\regle{redundant-set-state}) lisent le store ; les règles de deps lisent la
vue inter-rendus.
\end{invariant}

Pour \code{ObjChurn}, la certitude vient d'ailleurs : de la \emph{structure}.
La dep \code{obj} \emph{est} le slot (ligne \code{exact} de
\relation{effect_triggers}, \cref{sec:react-modele-effets}), l'effet écrit dans
ce slot une valeur \code{PerRender} sur tous ses chemins, et React ré-exécute
un effet dont une dep a changé. Trois faits must, donc une \Error. Le bras
correspondant d'\regle{infinite-loop} et le graphe de churn qui le généralise
sont l'objet du \cref{chap:churn}.

\begin{decision}[ADR-017 --- stabilité versionnée]
\emph{Contexte.} Cinq faux positifs \regle{always-unstable-deps} du premier
corpus avaient tous la forme de \code{ObjStateDep} : un état objet lu comme
« neuf à chaque rendu ». Sous ce FP se cachait un FN : \code{ObjChurn} n'était
pas détecté, parce que le join de deux références fraîches ne croît pas.
\emph{Décision.} Un treillis à six points, fragment d'un produit may $\times$
must ; la conversion côté lecture dans \code{StateVal} ; un bras « churn »
dans \regle{infinite-loop}. \emph{Alternatives refusées.} Rendre les lectures
d'état stables sans nouveau signal (rouvre le FN) ; encoder un compteur
d'événements dans le store pour faire tirer le widening (« a non-standard
hack ») ; garder dans le treillis deux points « change à coup sûr au set »,
qui n'avaient qu'un consommateur, lequel a de toute façon besoin d'un
raisonnement d'atteignabilité et vit donc dans la règle.
\end{decision}

\paragraph{Les mémos.} Le moteur ne calcule pas la \emph{valeur} d'un mémo,
seulement son identité, comme join des stabilités de ses deps :

\rustsnippet{src/domains/transfer/state_value.rs}{63}{77}{\code{recompute_memo} : les cas sans deps lisibles et \code{[]}}

Sans tableau lisible, le mémo peut être recalculé à chaque rendu comme ne
jamais l'être : $\Top$. Avec \code{[]} (d'arité exactement nulle), il est
figé : \code{Stable}. Sinon (\src{src/domains/transfer/state_value.rs}{78}{100}),
chaque dep contribue : une dep qui \emph{est} un slot contribue
\code{Versioned({slot})}, les autres leur stabilité évaluée. Pour
\code{Search}, \code{stable} hérite ainsi de la stabilité de la prop
\code{query}, qui n'est pas prouvée \code{PerRender} : silence.

\begin{definition}{Soundness en couches}{soundness-couches}
Un argument de soundness est \terme{en couches} lorsqu'il repose sur une
hypothèse du modèle abstrait dont chaque violation est signalée par une règle
dédiée. La conversion côté lecture suppose que les sets ont lieu hors du rendu :
si un setter est appelé \emph{pendant} le rendu avec une référence fraîche, le
slot change réellement à chaque rendu et \code{Versioned} le sous-estime. Cette
violation est signalée à part par \regle{setter-in-render}. De même, lire un
champ d'un objet d'état suppose qu'on ne le mute pas en place, violation
signalée par \regle{state-mutation}.
\end{definition}

\begin{piege}
\code{Versioned} n'est pas \code{Stable}. Un état jamais écrit se lit
\code{Versioned}, pas \code{Stable} (\issue{41}) : une dep sur lui ne peut pas
être omise, ce qui est aussi la position d'eslint. Et le fait qu'une dep
\code{Versioned} « garde » un effet dans \regle{infinite-loop} n'est sound
qu'\emph{avec} le bras churn : retirer l'un sans l'autre rouvre le faux négatif
d'\code{ObjChurn}. L'ADR-017 appelle ce couplage « a load-bearing soundness
dependency ».
\end{piege}

\section{Les règles des hooks et l'identité positionnelle}
\label{sec:react-modele-regles-hooks}

\subsection{Un hook est identifié par sa position}

Rien, dans \code{const [s] = useState(0)}, ne dit à React \emph{quel} état
rendre : ni nom, ni clé. React s'en remet à l'ordre des appels. Au $k$-ième
appel de hook d'un rendu, il rend le $k$-ième état de sa liste. Cet ordre doit
donc être le même à chaque rendu.

\begin{react}[règles des hooks]
\og Don't call Hooks inside loops, conditions, nested functions, or
\code{try}/\code{catch}/\code{finally} blocks. Instead, always use Hooks at the
top level of your React function, before any early returns.\fg{}
\cite{ReactRulesOfHooks} React-tRace donne la raison : \og each Hook call is
identified by its position in a linked list of internal bookkeeping
objects\fg{} \cite{LeeAhnYi2025}.
\end{react}

\begin{exemple}{\code{Cond}}{canon-cond}
\begin{lstlisting}[language=TypeScript,caption={\code{Cond}, transcrit de React-tRace (\file{/tmp/ch02/cond.tsx})},label={lst:react-modele-cond}]
import { useState } from "react";

export function Cond() {
  const [b, setB] = useState(false);
  if (b) {
    const [s] = useState(0);
    return <p>{s}</p>;
  }
  return <button onClick={() => setB((v) => !v)}>Show</button>;
}
\end{lstlisting}
\begin{sortie}
  Cond  (3 hooks)  cond.tsx
    error  conditional-hook  [hook:1]  (line 6:10)  this hook is called conditionally (not on every render path)
       → guarded by a condition evaluated here, so some render paths skip the hook (line 5:6)
\end{sortie}
\end{exemple}

Au premier rendu, \code{b} est faux, un seul hook est appelé. Après le clic,
\code{b} est vrai et un deuxième \code{useState} apparaît, que React n'a jamais
vu ; dans les mots de l'article, \og instead of displaying the initial state of
s which is 0, the UI breaks down at runtime\fg{}. L'analyseur le signale en
\Error{} : le défaut se produit dès que le code s'exécute sur l'autre branche,
et l'appel conditionnel est prouvé.

\subsection{Ce que fait le code}

Le label d'un hook est un entier (\srcline{src/ir/types.rs}{2} :
\code{pub type HookLabel = usize;}), attribué par le lowering dans l'ordre des
appels du corps. C'est l'exact analogue du « $\ell$ » que React-tRace attache
à chaque \code{useState} ; les handlers extraits partagent cet espace de
labels (d'où \code{[hook:1]} pour le \code{useState} conditionnel de
\code{Cond}, et « 3 hooks » en comptant le handler du bouton). L'IR range
chaque hook dans une variante de \code{HookEntry}, dont voici les deux
premières :

\rustsnippet{src/ir/hooks.rs}{247}{259}{les deux premières variantes de \code{HookEntry}}

\begin{table}[htbp]\centering\small
\begin{tabularx}{\linewidth}{@{}l>{\raggedright\arraybackslash}Xl@{}}
  \toprule
  Variante & Hooks React correspondants & Champs propres \\
  \midrule
  \code{State} & \code{useState}, \code{useReducer} (réducteur jeté) & \code{init} \\
  \code{Effect} & \code{useEffect}, \code{useLayoutEffect}, \code{useInsertionEffect} & \code{body_cfg}, \code{deps} \\
  \code{Memo} & \code{useMemo} & \code{body_cfg}, \code{deps} \\
  \code{Callback} & \code{useCallback} & \code{body_cfg}, \code{params}, \code{deps} \\
  \code{Ref} & \code{useRef} & \code{init} \\
  \code{Custom} & tout autre \code{use*} : hook du projet, de bibliothèque, ou hook React non modélisé (\code{useContext}\dots) & \code{name}, \code{args}, \code{binding}, provenance \\
  \code{Handler} & aucun : handler JSX \code{onX} ou callback d'\code{addEventListener} & \code{event}, \code{body_cfg} \\
  \bottomrule
\end{tabularx}
\caption{Les variantes de \code{HookEntry} (\src{src/ir/hooks.rs}{247}{313}) et
les hooks React qu'elles représentent. Toutes portent un \code{label} et un
\code{span} ; le classement est fait par \code{make_hook_entry}
(\src{src/lowering/hook_extractor.rs}{642}{747}).}
\label{tab:react-modele-hookentry}
\end{table}

La \cref{tab:react-modele-hookentry} donne la carte complète ; le détail des
champs est au \cref{chap:ir}. La règle \regle{conditional-hook} traduit la
règle de React en termes de CFG :

\rustsnippet{src/rules/impls/conditional_hook.rs}{3}{8}{\regle{conditional-hook} : un hook conditionnel ne domine pas toutes les sorties}

« Appelé à chaque rendu » devient « le bloc du hook domine toutes les sorties du
rendu » : tout chemin de l'entrée à un \code{return} passe par lui. Dans
\code{Cond}, le bloc \code{then} ne domine pas le \code{return} du bouton. La
dominance est définie au \cref{chap:interpretation-abstraite}
(\cref{def:dominance}) ; c'est un fait sur \emph{tous} les chemins,
donc un fait must, qui peut porter une \Error.

\begin{decision}[ADR-025 --- la fin d'un corps est un \code{return}]
Un corps qui « tombe » à la fin sans \code{return} renvoie \code{undefined} en
JavaScript. Le lowering scellait autrefois ce bloc par \code{Unreachable}, ce
qui, après inlining, coupait certains composants de leur propre sortie : 208
composants du corpus avaient ainsi un CFG de rendu « sectionné », et la
dominance des sorties se calculait sur moins de chemins que le programme n'en a.
\emph{Décision} (\adr{25}) : la fin de corps est \code{Return(undefined)} ;
\code{Unreachable} ne veut plus dire que « le contrôle s'arrête » (un
\code{throw}) ; un \code{return} inatteignable n'est pas une sortie.
\emph{Alternative refusée} : raccorder tout \code{Unreachable} d'un callee à la
jonction, « just an over-approximation ». Mesurée, elle inventait un chemin
qui atteint la sortie sans passer par les hooks du callee et signalait en
\Error{} l'idiome conforme \og guard-throw\fg{}
(\code{if (ctx === undefined) throw ...} puis \code{useState}). En React, un
\code{throw} pendant le rendu \emph{abandonne} le rendu : aucun chemin ne
continue vers les hooks suivants.
\end{decision}

\subsection{Appeler un setter pendant le rendu}

La seconde règle de pureté concerne le setter. L'appeler pendant le rendu est
autorisé, mais à une condition précise.

\begin{react}[set pendant le rendu]
\og Calling the set function \emph{during rendering} is only allowed from within
the currently rendering component. React will discard its output and
immediately attempt to render it again with the new state.\fg{} Un appel
inconditionnel donne l'erreur \og Too many re-renders. React limits the number
of renders to prevent an infinite loop.\fg{} \cite{ReactDocs} React-tRace
(§3.1.2) : \og an unconditional top-level call to a setter function causes an
infinite loop and is never correct\fg{}, même quand la valeur écrite est
toujours la même ; le runtime lève une exception après 25 tentatives.
\end{react}

\begin{exemple}{\code{Inf2}}{canon-inf2}
\begin{lstlisting}[language=TypeScript,caption={\code{Inf2} : l'updater identité, appelé dans le corps (\file{/tmp/ch02/setter_in_render.tsx})},label={lst:react-modele-inf2}]
import { useState } from "react";

export function Inf2() {
  const [s, setS] = useState(0);
  setS((x) => x);
  return <div>{s}</div>;
}
\end{lstlisting}
\begin{sortie}
  Inf2  (1 hooks)  setter_in_render.tsx
    error  setter-in-render  [hook:0]  (line 5:2)  setter `setS` called directly in the render body, move this call into a useEffect or an event handler
       → `setS` is a state setter, so calling it writes state (line 5:2)
\end{sortie}
\end{exemple}

L'appel est dans le bloc d'entrée, qui domine toutes les sorties : \Error.
L'exemple est instructif pour le bail-out : l'updater identité ne change rien,
et pourtant React boucle, parce qu'un set pendant le rendu fait \emph{jeter} le
rendu avant toute comparaison (\cref{sec:react-modele-trace}).

La documentation de React prévoit pourtant un usage légitime : \og ajuster
l'état pendant le rendu\fg{} quand une prop change, à la place d'un effet
(page \og You Might Not Need an Effect\fg{}).

\begin{lstlisting}[language=TypeScript,caption={\code{List} : le motif « adjust state during render » (\file{/tmp/ch02/adjust.tsx})},label={lst:react-modele-adjust}]
import { useState } from "react";

export function List({ items }: { items: string[] }) {
  const [prevItems, setPrevItems] = useState(items);
  const [selection, setSelection] = useState<string | null>(null);
  if (items !== prevItems) {
    setPrevItems(items);
    setSelection(null);
  }
  return <ul onClick={() => setSelection(items[0])}>{selection}</ul>;
}
\end{lstlisting}

\begin{sortie}
  List  (3 hooks)  adjust.tsx
    warn   setter-in-render  [hook:1]  (line 8:4)  setter `setSelection` called directly in the render body, move this call into a useEffect or an event handler
\end{sortie}

La règle fait trois choses différentes, décrites par sa documentation :

\rustsnippet{src/rules/impls/setter_in_render.rs}{18}{36}{les niveaux de \regle{setter-in-render}}

\code{setPrevItems(items)} est tu : la garde qui le domine (\code{items !==
prevItems}) lit le slot qu'il écrit, et devient fausse dès que la valeur écrite
y est (\code{converges_once_written}, \src{src/engine/guards.rs}{60}{67}).
\code{setSelection(null)} est signalé en \Warning{} alors qu'il fait partie du
même motif documenté : sa garde ne lit pas \code{selection}, et la règle ne
peut pas prouver que l'appel converge. C'est un faux positif de niveau
\Warning, que la doctrine tolère.

\section{Strict Mode : le double rendu}
\label{sec:react-modele-strictmode}

En développement, React peut exécuter volontairement deux fois une partie du
code pour faire apparaître les impuretés.

\begin{react}[\code{<StrictMode>}]
Strict Mode active des comportements \emph{de développement} : \og Your
components will re-render an extra time to find bugs caused by impure
rendering\fg{}, \og Your components will re-run Effects an extra time to find
bugs caused by missing Effect cleanup\fg{}, et de même pour les ref callbacks.
Sont appelés deux fois : \og Your component function body (only top-level
logic, so this doesn't include code inside event handlers)\fg{} et les fonctions
passées à \code{useState}, aux fonctions \code{set}, à \code{useMemo} et à
\code{useReducer}. Pour les effets, \og React will run one extra
development-only setup+cleanup cycle before the first real setup\fg{}.
\og All of these checks are development-only and do not impact the production
build.\fg{} \cite{ReactDocs}
\end{react}

Un composant impur montre la différence. Celui-ci compte ses rendus dans une
variable de module :

\begin{lstlisting}[language=TypeScript,caption={\code{Impure} : le rendu modifie une variable de module (\file{/tmp/ch02/strict.tsx})},label={lst:react-modele-impure}]
import { useState } from "react";

let renders = 0;

export function Impure() {
  const [n, setN] = useState(0);
  renders = renders + 1;
  return <button onClick={() => setN(n + 1)}>{renders}</button>;
}
\end{lstlisting}

Sous Strict Mode, le compteur avance de deux par rendu en développement et de
un en production : c'est exactement ce que le double appel veut révéler.
\reactant{} ne dit rien de ce composant :

\begin{sortie}
  Impure  (2 hooks)  strict.tsx  ✓
    verified  conditional-hook  all hooks run unconditionally, in a stable order
    verified  lazy-init  no useState/useRef initializer re-runs work on every render
    verified  setter-in-render  no setter is called during render
    verified  state-mutation  no state or prop object is mutated in place
\end{sortie}

Aucune règle du catalogue ne porte sur l'écriture d'une variable de module
pendant le rendu ; \regle{state-mutation} ne vise que les objets d'état et de
props, comme sa ligne \code{verified} le dit.

\paragraph{Ce que fait le code.} Rien : Strict Mode n'est pas modélisé par le
moteur. La seule trace est documentaire, dans \regle{missing-cleanup}, qui
rappelle qu'\og under StrictMode it runs it twice on the first mount\fg{}
(\src{src/rules/impls/missing_cleanup.rs}{9}{13}) et reprend le fait dans son
message (\cref{lst:react-modele-listener}).

Ce silence est-il sound ? L'argument suivant n'est écrit nulle part dans le
code ; c'est une déduction. Le point fixe calcule un ensemble d'états
atteignables et joint les écritures. Exécuter deux fois un rendu, un
initialiseur ou un updater \emph{purs} produit deux fois les mêmes écritures, et
$x \join x = x$ : l'ensemble calculé est inchangé. Pour les effets, le cycle
supplémentaire setup--cleanup n'ajoute que des exécutions d'un corps que le
point fixe exécute déjà à chaque itération. Strict Mode ne rend donc visible
que ce qui viole la pureté, c'est-à-dire des hypothèses que le modèle abstrait
fait déjà (\cref{sec:react-modele-rendu}).

\begin{piege}
Le double rendu de Strict Mode n'est pas un rendu « de trop » au sens des
règles. Un contributeur qui voudrait le modéliser comme deux tours ferait
tirer \regle{unnecessary-rerender} ou compterait deux fois une écriture de
montage ; or ces règles parlent de la production, où le double appel n'existe
pas. Inversement, un utilisateur qui voit son effet de montage s'exécuter deux
fois en développement ne doit pas en conclure qu'il boucle, ni que
\reactant{} se trompe en le traitant comme un effet de montage : l'effet est
exécuté une fois de plus \emph{par construction}, et la règle qui s'en préoccupe est
\regle{missing-cleanup}, pas \regle{infinite-loop}.
\end{piege}

\section{Le contexte et la propagation du rendu}
\label{sec:react-modele-contexte}

\subsection{Rendre un parent, c'est rendre ses enfants}

Un rendu est récursif : rendre un composant rend les composants qu'il renvoie
(\cref{sec:react-modele-rendu}). Un set dans un parent re-rend donc tout son
sous-arbre, sauf les parties que React peut prouver inchangées (un élément
reçu en \code{children} dont l'identité est conservée, un composant enveloppé
par \code{memo} dont les props sont égales). Le contexte est l'autre canal de
propagation : il traverse l'arbre sans passer par les props.

\begin{react}[\code{useContext}]
\og useContext returns the context value for the calling component. It is
determined as the value passed to the closest SomeContext above the calling
component in the tree.\fg{} À défaut de provider, c'est la \code{defaultValue}
de \code{createContext}. \og React automatically re-renders all the children
that use a particular context starting from the provider that receives a
different value. The previous and the next values are compared with the
\code{Object.is} comparison. Skipping re-renders with memo does not prevent the
children receiving fresh context values.\fg{} D'où le conseil d'envelopper la
valeur fournie dans \code{useMemo} et les fonctions dans \code{useCallback}
\cite{ReactDocs}.
\end{react}

\begin{lstlisting}[language=TypeScript,caption={Un provider dont la valeur est neuve à chaque rendu (\file{/tmp/ch02/context.tsx})},label={lst:react-modele-context}]
import { createContext, useContext, useState } from "react";

const ThemeCtx = createContext({ dark: false, toggle: () => {} });

export function ThemeProvider({ children }: { children: any }) {
  const [dark, setDark] = useState(false);
  return (
    <ThemeCtx.Provider value={{ dark, toggle: () => setDark((d) => !d) }}>
      {children}
    </ThemeCtx.Provider>
  );
}

export function Button() {
  const { dark } = useContext(ThemeCtx);
  return <button>{dark ? "dark" : "light"}</button>;
}
\end{lstlisting}

\begin{sortie}
  Button  (1 hooks)  context.tsx
    info   analysis-limit  [hook:0]  (line 15:8)  hook `useContext` was not found in the registry. Pass its source file or add a HookSummary to analyse it (FN possible)
    suspended  analysis-limit  1 passing check(s) withheld: the analysis was truncated in this component, so they are not guaranteed
  ThemeProvider  (1 hooks)  context.tsx
    info   analysis-limit  component `ThemeCtx.Provider` was not found in the analysis registry. Pass its file on the command line to analyse it (FN possible)
    warn   unstable-context-value  (line 8:4)  `ThemeCtx.Provider` is given a newly allocated value on every render. `Object.is` fails for every consumer, so each `useContext(ThemeCtx)` re-renders whenever this component does, even when nothing in the value changed; wrap the value in `useMemo`
    suspended  analysis-limit  4 passing check(s) withheld: the analysis was truncated in this component, so they are not guaranteed
\end{sortie}

Les deux côtés du contexte sont traités très différemment.

\paragraph{Côté consommateur : $\Top$.} \code{useContext} n'est pas modélisé.
Le lowering le range, comme tout hook React sans modèle, dans une entrée
\code{Custom} dont la provenance garde le spécificateur \code{react}
(commentaire de \src{src/lowering/hook_extractor.rs}{730}{732}) :

\rustsnippet{src/lowering/hook_extractor.rs}{733}{744}{un hook React non modélisé devient une entrée \code{Custom}}

Son site d'appel est un marqueur \code{HookMarker(_, Unknown)}, qui s'évalue en
$\Top$ ; c'est la l.~143 de \file{src/domains/transfer/state_value.rs}, déjà
citée au \cref{sec:react-modele-identite}. Le consommateur lit
donc « n'importe quelle valeur », ce qui est sound, et l'\Info{}
\regle{analysis-limit} rend le trou visible ; les vérifications qui auraient
réussi sont « suspendues » plutôt qu'affichées comme prouvées. C'est, d'après
\file{docs/limitations.md}, la première source de limites d'analyse du corpus
(363 sites, \issue{28}).

\paragraph{Côté provider : l'identité.} La valeur passée à \code{value} est un
littéral objet, donc \code{PerRender} : fait must. La règle
\regle{unstable-context-value} en tire un \Warning{} :

\rustsnippet{src/rules/impls/unstable_context_value.rs}{4}{13}{le défaut visé par \regle{unstable-context-value}}

\Warning{} et non \Error{} : la référence fraîche est certaine, mais son coût
(des rendus de consommateurs) ne l'est pas --- c'est la seconde moitié de la
définition du \Warning{} (« fait certain au coût incertain »). Au-delà de cette
règle, la relation \relation{context_consumers} (\adr{32}) et la dépendance de
rendu (\adr{41}) suivent un contexte prouvé de son provider jusqu'aux
consommateurs placés sous lui, pour les règles de cascade
\regle{state-lifted-too-high} et \regle{wasted-subtree-render}
(\cref{chap:inter-composants}). Un élément que l'analyse ne peut pas ouvrir, en
particulier un composant enveloppé par \code{memo} ou \code{forwardRef}
(\issue{64}), y est compté comme un usage possible.

\begin{remarque}
La première ligne \code{info} de \code{ThemeProvider} signale que
\code{ThemeCtx.Provider} n'est pas un composant du registre : l'analyseur voit
une balise JSX dont il ne connaît pas le corps, et le dit. Le provider n'en
est pas moins reconnu par la règle, qui repère \code{ThemeCtx} comme un
\code{createContext(...)} de niveau module.
\end{remarque}

\section{Server Components et la directive \code{"use client"}}
\label{sec:react-modele-rsc}

Tout ce qui précède suppose que le composant vit dans le navigateur, où il a un
état et des effets. Les Server Components (RSC), popularisés par l'App Router
de Next.js, rompent cette hypothèse.

\begin{react}[Server et Client Components]
Un Server Component est rendu une fois, sur le serveur, et ne persiste pas :
\og Server Components cannot use most Hooks. When Server Components are
rendered, their output is essentially a list of components for the client to
render. Server Components do not persist in memory after render and cannot
have their own state.\fg{} La frontière est une directive de module :
\og Add \code{'use client'} at the top of a file to mark the module and its
transitive dependencies as client code\fg{}, et elle \og must be at the very
beginning of a file, above any imports or other code\fg{}. Il n'y a pas de
directive pour le côté serveur, qui est le défaut ; et \og a single module may
be evaluated on the server when imported from server code and on the client
when imported from client code\fg{} \cite{ReactDocs}.
\end{react}

La dernière phrase est la difficulté : être « serveur » n'est pas une propriété
d'un fichier, mais de la façon dont on l'atteint dans le graphe d'imports.
Considérons un projet Next minimal (\file{/tmp/ch02next/}) : un
\file{next.config.js} vide, une page \file{app/page.tsx} sans directive, et un
composant \file{components/like-button.tsx} qui commence par
\code{"use client"}.

\begin{lstlisting}[language=TypeScript,caption={\file{app/page.tsx} : un état dans un Server Component (\file{/tmp/ch02next/app/page.tsx})},label={lst:react-modele-page}]
import { useState } from "react";
import { LikeButton } from "../components/like-button";

export default function Page() {
  const [open, setOpen] = useState(false);
  return (
    <main onClick={() => setOpen(!open)}>
      <LikeButton />
    </main>
  );
}
\end{lstlisting}

\begin{shell}
$ cd /tmp/ch02next && reactant check --no-color --info --show-clean .
\end{shell}
\begin{sortie}
[warn] no tsconfig `paths` found, so aliased imports (e.g. `@/...`) stay unresolved and their targets are NOT analyzed (possible false negatives). Aliases declared only in next.config are not read.
  LikeButton  (2 hooks)  components/like-button.tsx  ✓
    verified  conditional-hook  all hooks run unconditionally, in a stable order
    verified  lazy-init  no useState/useRef initializer re-runs work on every render
    verified  setter-in-render  no setter is called during render
    verified  state-mutation  no state or prop object is mutated in place
  Page  (2 hooks)  app/page.tsx
    warn   server-component-hook  [hook:0]  (line 5:8)  `useState` is called in a Server Component. this file is an App Router `page` and no `"use client"` directive covers it, so React renders it on the server, where hooks do not exist; add `"use client"` at the top of the file, or move the stateful part into a child component that declares it
\end{sortie}

Le projet est reconnu comme Next.js par son \file{next.config.js}. La page est
une entrée de l'App Router, aucune directive ne la couvre : c'est un module
serveur, et son \code{useState} est un défaut. \code{LikeButton}, sous
\code{"use client"}, est un composant client ordinaire.

\begin{decision}[ADR-026 --- analyser les Server Components quand même]
\emph{Décision} (\adr{26}). Chaque module garde ses faits bruts
(\code{ModuleFacts {directives, imports}}) ; un module est serveur s'il est
atteignable depuis une entrée de l'App Router (\code{page}, \code{layout},
\code{template}, \code{default}, \code{not-found}, \code{loading}) sans
franchir de \code{"use client"} ; la règle \regle{server-component-hook} y
signale les hooks, au niveau \Warning{} ; elle se tait si le programme
n'utilise jamais \code{"use client"}. \emph{Alternatives refusées.}
\emph{Sauter} les modules serveur : décider « serveur, on passe » sur un graphe
d'imports incomplet ferait disparaître leurs défauts, un faux négatif. Une
preuve must (\code{Certified}) pour la règle : \og would dress a path
convention as a domain proof\fg{}, d'où le plafond \Warning{} : les deux faits porteurs (une convention de noms de fichiers, un graphe dont le résolveur a pu manquer des arêtes) sont hors du domaine abstrait. Un seul diagnostic par composant, car le défaut est la directive manquante du module. Supprimer les
autres diagnostics dans un module serveur : toute erreur de classification
deviendrait un faux négatif. Analyser un Server Component ne coûte rien en
soundness : il n'a pas d'état, et l'interprétation abstraite le traite comme
un composant dont les props sont $\Top$.
\end{decision}

La règle sous-déclare quand un import sur le chemin d'un module serveur n'est
pas résolu (\issue{29}), et le premier avertissement de la sortie rappelle
justement que des imports aliasés peuvent rester non résolus.

\section{React-tRace : une sémantique formelle des hooks}
\label{sec:react-modele-trace}

Les sections précédentes décrivent React à partir de sa documentation, qui est
informelle : elle ne dit pas, par exemple, \emph{quand exactement} la file
d'updaters est rejouée (« sometime during the next render », relèvent les
auteurs). Un analyseur par interprétation abstraite a besoin d'une sémantique
concrète précise à sur-approximer. React-tRace, de Jay Lee, Joongwon Ahn et
Kwangkeun Yi (OOPSLA 2025) \cite{LeeAhnYi2025}, en fournit une pour le cœur
des hooks : \code{useState} et \code{useEffect}. Cette section en résume les
éléments dont le projet se sert ; les notations sont celles de l'article.

\subsection{Le langage}

L'article travaille sur un petit langage applicatif, découplé de JavaScript
(\og we decouple ourselves from the subtleties of the complex semantics of
JS\fg{}). Sa syntaxe (Fig.~1 de l'article) :
\[
\begin{array}{r@{\;}c@{\;}l}
  P &::=& D\;e \qquad D ::= \mathtt{let}\ C(x) = e\\[2pt]
  e &::=& () \mid \mathtt{true} \mid \mathtt{false} \mid n \mid x \mid C
          \mid e \oplus e \mid [e] \mid \mathtt{print}\ e
          \mid \mathtt{if}\ e\ \mathtt{then}\ e\ \mathtt{else}\ e \mid e;e\\
    &\mid& \mathtt{fun}\ x \to e \mid e\ e \mid \mathtt{let}\ x = e\ \mathtt{in}\ e
     \mid \mathtt{let}\ (x, x_{\mathit{set}}) = \mathtt{useState}_\ell\ e\ \mathtt{in}\ e
     \mid \mathtt{useEffect}\ e
\end{array}
\]
Un programme est une suite de définitions de composants suivie d'une
expression principale. \code{[e]} est une vue composée (l'analogue du JSX
imbriqué) ; une fermeture utilisée comme vue représente un handler (un bouton
\emph{est} son handler). Les hooks ne peuvent apparaître qu'au niveau supérieur
d'un corps de composant, ce que l'interpréteur vérifie à l'analyse syntaxique ;
chaque \code{useState} porte un label unique $\ell$. L'effet n'a pas de deps :
il s'exécute après chaque rendu, les deps étant \og a straightforward
extension\fg{}.

\subsection{Les objets sémantiques}

Les valeurs de base sont les constantes, les fermetures
$\langle \lambda x.e, \sigma\rangle$ et un tampon de sortie $\omega$. La
machinerie React ajoute :
\begin{itemize}
  \item la \emph{spécification de composant} $\langle C, v\rangle$ (un nom et
    son argument : appliquer un composant ne fait que l'empaqueter) ;
  \item le \emph{setter} $\langle \ell, p\rangle$ : le label du \code{useState}
    d'origine et le chemin $p$ de la vue qui le possède ;
  \item la \terme{mémoire d'arbre} $m = [p \mapsto \pi]$, qui associe à chaque
    chemin une \emph{vue}
    $\pi = \{\mathit{spec}, \mathit{dec}, \mathit{sttst}, \mathit{effq}, \mathit{child}\}$ :
    sa spécification, un ensemble de \emph{décisions}
    $d \in \{\mathsf{Check}, \mathsf{Effect}\}$, un store d'état
    $\rho = [\ell \mapsto \{\mathit{val} : v, \mathit{sttq} : q\}]$ (valeur
    courante et file d'updaters), une file d'effets et un arbre d'enfants ;
  \item une \emph{phase} d'évaluation $\phi \in \{\mathsf{Init}, \mathsf{Succ},
    \mathsf{Normal}\}$ : premier appel d'un corps, appels suivants, ou
    expression principale, effets et handlers ;
  \item un \emph{mode} du moteur : \emph{rendered} (l'interface est à l'écran,
    des effets attendent), \emph{check} (il faut vérifier si un rendu est
    nécessaire) ou \emph{event loop} (on attend une entrée).
\end{itemize}
Les décisions sont les deux mots de la citation du
\cref{sec:react-modele-effets} : \textsf{Check} marque une vue à re-vérifier
(un set l'a visée), \textsf{Effect} une vue dont les effets doivent tourner
après ce rendu.

\subsection{La boucle de rendu}

Le moteur est un système de transitions à quatre règles (Fig.~4 de
l'article). En omettant le tampon de sortie et la table des définitions
$\delta$, qui ne change pas :
\[
\begin{array}{l@{\qquad}l}
\textsc{StepInit} &
  \dfrac{[\,],[\,] \vdash e \Downarrow^{\mathsf{Normal}} s
         \qquad [\,] \vdash \mathit{init}(s) = \langle t, m\rangle}
        {\langle e\rangle \hookrightarrow \langle t, m, \text{rendered}\rangle}\\[14pt]
\textsc{StepEffect} &
  \dfrac{m \vdash \mathit{commitEffs}(t) = m'}
        {\langle t, m, \text{rendered}\rangle \hookrightarrow \langle t, m', \text{check}\rangle}\\[14pt]
\textsc{StepCheck} &
  \dfrac{m \vdash \mathit{check}(t) = \langle \mu, m'\rangle}
        {\langle t, m, \text{check}\rangle \hookrightarrow \langle t, m', \mu\rangle}\\[14pt]
\textsc{StepEvent} &
  \dfrac{\langle \lambda x.e, \sigma\rangle \in \mathit{handlers}(m, t)
         \qquad m, \sigma[x \mapsto ()] \vdash e \Downarrow^{\mathsf{Normal}} v, m'}
        {\langle t, m, \text{event loop}\rangle \hookrightarrow \langle t, m', \text{check}\rangle}
\end{array}
\]
avec $\mu \in \{\text{rendered}, \text{event loop}\}$. La
\cref{fig:react-modele-trace-modes} en donne l'automate. On y lit la
\cref{fig:react-modele-chronologie} sous forme exacte : \textsc{StepInit} est le
montage ; \textsc{StepEffect} exécute les effets en file ; \textsc{StepCheck}
re-rend les vues marquées \textsf{Check} et revient en mode \emph{rendered} si
un état a changé, en \emph{event loop} sinon ; \textsc{StepEvent} exécute un
handler, ce qui impose une nouvelle vérification.

\ifhastikz
\begin{figure}[htbp]\centering
\begin{tikzpicture}[->,>={Stealth[length=2mm]},node distance=30mm,
    every state/.style={draw=ra-blue,thick,fill=ra-bg,minimum size=17mm,font=\small,align=center},
    lab/.style={font=\scriptsize}]
  \node[state,initial,initial text={$\langle e, \delta\rangle$}] (r) {rendered};
  \node[state,right=of r] (c) {check\\{\scriptsize $\circlearrowleft$}};
  \node[state,right=of c] (e) {event\\loop};
  \draw (r) edge[bend left=20] node[lab,above]{\textsc{StepEffect}} (c);
  \draw (c) edge[bend left=20] node[lab,below,align=center]{\textsc{StepCheck}\\un état a changé} (r);
  \draw (c) edge[bend left=20] node[lab,above,align=center]{\textsc{StepCheck}\\rien n'a changé} (e);
  \draw (e) edge[bend left=20] node[lab,below]{\textsc{StepEvent}} (c);
\end{tikzpicture}
\caption{Les modes de React-tRace et ses quatre transitions. La flèche
d'entrée est \textsc{StepInit} (évaluation de l'expression principale et
initialisation de l'arbre). La boucle infinie d'effets est le cycle
rendered $\to$ check $\to$ rendered qui ne passe jamais par \emph{event loop} ;
un handler, lui, n'est exécuté qu'en mode \emph{event loop}.}
\label{fig:react-modele-trace-modes}
\end{figure}
\fi

\subsection{Les règles d'évaluation}

Sous la boucle, l'évaluation d'une expression
$\Sigma, \sigma \vdash e \Downarrow^{\phi}_{p} v, \Sigma', \omega$ est une
sémantique à grands pas. Les règles propres à React sont les suivantes
(Fig.~5 à 10 de l'article).
\begin{description}
  \item[\textsc{SttBind}] (phase \textsf{Init}) : \code{useState} évalue son
    argument $v_1$, crée l'entrée $\ell \mapsto \{\mathit{val}: v_1,
    \mathit{sttq}: [\,]\}$ et lie $x_{\mathit{set}}$ au setter
    $\langle \ell, p\rangle$.
  \item[\textsc{SttReBind}] (phase \textsf{Succ}) : \code{useState} rejoue la
    file $[\langle \lambda x_i.e_i, \sigma_i\rangle]_{i=1}^{n}$ dans l'ordre,
    $v_0 \to v_1 \to \dots \to v_n$, vide la file, stocke $v_n$, et ajoute la
    décision \textsf{Effect} si $v_n \not\equiv v_0$.
  \item[\textsc{AppSetComp}] (set pendant le corps, phases \textsf{Init} et
    \textsf{Succ}) : le chemin du setter doit être celui de la vue courante ;
    l'updater est ajouté à la file et la décision \textsf{Check} posée.
  \item[\textsc{AppSetNormal}] (set dans un effet ou un handler, phase
    \textsf{Normal}) : l'updater est ajouté à la file de la vue $m[p]$ désignée
    par le setter, quelle qu'elle soit, et \textsf{Check} y est posée. C'est la
    seule règle qui voie toute la mémoire : un enfant peut mettre à jour son
    parent.
  \item[\textsc{Eff}] : \code{useEffect e} ajoute la fermeture
    $\langle \lambda\_.e, \sigma\rangle$ à la file d'effets de la vue ---
    l'effet capture l'environnement \emph{du rendu}.
  \item[\textsc{EvalOnce} / \textsc{EvalMult}] : un corps est évalué, puis
    ré-évalué en phase \textsf{Succ} tant qu'il pose \textsf{Check} (file
    d'effets vidée à chaque essai) ; React lève une exception après 25
    essais.
  \item[\textsc{CommitEffsPath} / \textsc{CommitEffsPathIdle}] : les effets
    d'une vue ne sont exécutés que si sa décision contient \textsf{Effect},
    enfants d'abord (post-ordre), puis dans l'ordre syntaxique.
  \item[\textsc{CheckIdle} / \textsc{CheckNoEffect} / \textsc{CheckEffect}] :
    une vue sans \textsf{Check} est sautée ; avec \textsf{Check}, son corps
    est ré-évalué ; si la nouvelle décision ne contient pas \textsf{Effect}
    (l'état est revenu au même), rien n'est rendu ; sinon l'enfant est
    réconcilié et le mode redevient \emph{rendered}.
  \item[\textsc{ReconcileComEffect} / \textsc{ReconcileComNew}] : un enfant
    du même composant est ré-évalué en \textsf{Succ} ; un composant différent
    à la même place est initialisé à neuf.
\end{description}

\begin{exemple}{\code{SelfCounter} en React-tRace}{trace-selfcounter}
Au montage, \textsc{InitCom} évalue le corps en phase \textsf{Init}
(\textsc{SttBind} : $s = 0$ ; \textsc{Eff} met l'effet en file) et pose
\textsf{Effect} : mode \emph{rendered}. \textsc{StepEffect} exécute l'effet en
phase \textsf{Normal} : $s < 3$, donc \textsc{AppSetNormal} met
$\lambda s.\,s+1$ en file et pose \textsf{Check} ; mode \emph{check}.
\textsc{StepCheck} ré-évalue le corps : \textsc{SttReBind} rejoue la file,
$0 \to 1$, $1 \not\equiv 0$ donc \textsf{Effect} : \textsc{CheckEffect}, mode
\emph{rendered}. Le cycle se répète jusqu'à $s = 3$ ; alors l'effet tourne mais
ne pose pas \textsf{Check}, et \textsc{StepCheck} (\textsc{CheckIdle}) mène en
\emph{event loop}.
\end{exemple}

\subsection{Les théorèmes}

L'article prouve deux propriétés que la documentation de React énonce
informellement (§6.1).

\begin{theoreme}{Ré-évaluation par appel d'un setter (Theorem 1)}{trace-reevaluation}
La dérivation de l'évaluation retentée d'un corps contient plusieurs évaluations
de ce corps si et seulement si \textsc{AppSetComp} apparaît dans la dérivation
de la première évaluation.
\end{theoreme}

\begin{theoreme}{Condition d'exécution des effets (Theorem 2)}{trace-effets}
Si $\langle t, m, \omega, \delta, \text{check}\rangle \hookrightarrow
\langle t, m', \omega', \delta, \mu\rangle$, alors pour tout chemin $p$
atteignable depuis la racine dans $m$ et $m'$, \textsc{CommitEffsPath} en $p$
apparaît dans la dérivation de la transition suivante si et seulement s'il
existe un ancêtre $p'$ de $p$ tel que (1) pour un $\ell$ du store de $m'[p']$,
$m[p'].\mathit{sttst}[\ell].\mathit{val} \neq m'[p'].\mathit{sttst}[\ell].\mathit{val}$ ;
ou (2) la dérivation de la transition contient \textsc{AppSetComp} de chemin
$p'$.
\end{theoreme}

Le premier théorème est le mécanisme d'\code{Inf2} : un set dans le corps fait
jeter l'évaluation et recommencer, indépendamment de la valeur écrite. Le
second dit quand les effets tournent : après un changement d'état de la vue
\emph{ou d'un de ses ancêtres} (c'est la cascade de rendu), ou si un set a eu
lieu pendant le corps. Un troisième résultat (Theorem~8, appuyé sur le
Lemma~7) concerne une optimisation de React, qui applique par avance les
updaters en tête de file sans ré-évaluer le corps : si ces updaters sont purs
(Definition~3), l'état final est le même. Il justifie que la sémantique ignore
cette optimisation. La validité d'une vue (Definition~6) repose enfin sur le
fait que \og each useState call in a component body is executed exactly once
during initialization\fg{} --- c'est ici que la règle des hooks entre dans la
formalisation.

\paragraph{Conformité.} La suite de tests couvre les 44 règles d'évaluation par
38 tests répartis en 18 scénarios (Table~1), rejoués sur React 16.14.0,
17.0.2, 18.3.1 et 19.1.0 ; une seule différence, due à l'optimisation
ci-dessus (scénario S17, set depuis un événement utilisateur). L'interpréteur
et le visualiseur sont publiés (\code{Zeta611/react-trace}).

\paragraph{Extensions esquissées.} Le §7 de l'article esquisse
\code{useRef} (un store de refs sans \textsf{Check}), \code{useMemo} (valeur
et deps, recalcul en \textsf{Succ} si les deps diffèrent), \code{useContext}
(recherche ascendante du provider, \textsf{Check} des consommateurs quand la
valeur change), \code{useLayoutEffect} (un mode supplémentaire) et les hooks
personnalisés (de simples compositions). Aucune n'est formalisée : c'est
précisément la part que \reactant{} doit ajouter sans garantie formelle.

\section{La sémantique concrète du projet : ce qu'on garde, ce qu'on abstrait}
\label{sec:react-modele-concrete}

\subsection{Le contrat : ADR-001}

\begin{decision}[ADR-001 --- React-tRace comme sémantique concrète de référence]
\emph{Contexte.} Une interprétation abstraite dérive une sémantique abstraite
$C^\sharp$ d'une sémantique concrète $C$ ; sans $C$ explicite, la soundness ne
peut pas être établie et les fonctions de transfert \og are written by
guesswork\fg{}. \emph{Décision} (\adr{1}). React-tRace est la sémantique
concrète $C$ ; les fonctions de transfert en dérivent ; les extensions (deps,
\code{useMemo}, \code{useCallback}, \code{useRef}, objets) sont spécifiées
comme extensions de React-tRace. \emph{Justification.} C'est la seule
formalisation publique de React accompagnée d'une preuve de conformité ; la
boucle \textsc{StepInit} $\to$ \textsc{StepEffect} $\to$ \textsc{StepCheck}
correspond directement à l'itération de point fixe ; \textsc{SttReBind},
\textsc{CheckEffect} et \textsc{CheckNoEffect} définissent exactement les
conditions de re-rendu. \emph{Limites acceptées.} React-tRace ne couvre que
\code{useState} et \code{useEffect} sans deps ; son langage n'est pas
JavaScript ; les extensions n'ont pas de garantie formelle équivalente.
\end{decision}

L'ADR annonce aussi trois conséquences que le dépôt, au commit \code{e67b10a},
ne réalise pas. Le fichier \file{docs/semantics.md} censé spécifier les
extensions n'existe pas ; aucune fonction de \file{src/domains/} ne cite de
règle React-tRace (\code{grep -rn "tRace" src} ne trouve rien) ; aucun test ne
compare l'analyseur à l'interpréteur de l'article sur ses exemples (aucun test
ne mentionne \code{Flicker}, \code{SelfCounter} ou \code{Inf2}). React-tRace est
donc la référence \emph{conceptuelle} du projet --- l'ADR-004 et l'ADR-012 s'y
réfèrent, et ce chapitre en tire la table de correspondance ---, pas un oracle
outillé.

\subsection{De la boucle de rendu à la boucle de point fixe}

\begin{decision}[ADR-004 --- un CFG de rendu, un CFG par effet]
Chaque composant est représenté par un \code{render_cfg} et une liste de
\code{HookEntry}, dont chaque effet porte son propre \code{body_cfg}
(\adr{4}). L'alternative refusée est un méta-CFG unique, rendu et effets
reliés par des arêtes rendu $\to$ effet $\to$ check : un grand graphe
difficile à maintenir, qui rend implicite la séparation rendu/effet et
\og risks incorrect modeling of React's execution order\fg{}. Or la séparation
est un invariant de React (\og the effect never executes during the
render\fg{}), et les défauts visés sont des propriétés de l'\emph{état} au
point fixe, pas des chemins d'un graphe unifié.
\end{decision}

La boucle de \code{analyze_component_impl}
(\src{src/engine/fixpoint.rs}{130}{698}) joue abstraitement les tours de React.
L'\cref{alg:react-modele-fixpoint} en donne une lecture fidèle, réduite à ce
qui concerne ce chapitre ; son étude complète est au
\cref{chap:point-fixe-composant}.

\begin{algorithm}[htbp]
\Entree{un composant : \code{render_cfg}, \code{hooks}, environnement initial (props)}
\Sortie{les stores convergés, \code{widen_trace}, \code{effect_setter_writes}}
$S \leftarrow \Bot$\;
\PourCh{$\mathtt{State}\{\ell, \mathit{init}\}$}{$S[\ell] \leftarrow S[\ell] \join \sem{\mathit{init}}$ \tcp*{\textsc{SttBind}, l.~314--353}}
\Tq{vrai}{
  $(\mathit{env}, R) \leftarrow$ passe de rendu sur $S$ \tcp*{\textsc{StepInit}/\textsc{Succ}}
  recalculer les mémos depuis $\mathit{env}$\;
  $E \leftarrow \bigsqcup_{\text{tout effet}}$ passe d'effet depuis $\mathit{env}$ sur $S$ \tcp*{\textsc{StepEffect}, deps ignorées}
  $H \leftarrow \bigsqcup_{\text{tout handler}}$ passe de handler depuis $\mathit{env}$ sur $S$ \tcp*{\textsc{StepEvent}, 0..N fois}
  $N \leftarrow R \join E \join H \join$ écritures des autres composants\;
  \Si{$N \lleq S$}{\textbf{sortir} \tcp*{\textsc{StepCheck} : rien n'a changé}}
  $i \leftarrow i+1$\;
  \eSi{$i \geq 3$}{\code{widen_trace} $\cup$= slots changés dans $R \join E$ (pas $H$)\; $S \leftarrow S \widen N$ (à seuils)\;}{$S \leftarrow N$\;}
}
rejouer le rendu sur les stores convergés ; rejouer les effets depuis $\Bot$ $\to$ \code{effect_setter_writes}\;
\caption{La boucle de point fixe d'un composant (lecture de \src{src/engine/fixpoint.rs}{355}{593}, élaguée : garde-fou à 100 itérations, relations calculées à la fin)}
\label{alg:react-modele-fixpoint}
\end{algorithm}

Une itération abstraite n'est pas un tour de React : elle en représente
\emph{tous} à la fois. Le store $S$ de l'itération $k$ sur-approxime les
valeurs que l'état peut avoir après au plus $k$ tours, de n'importe quel type
(déclenché par un effet ou par un handler), et la boucle s'arrête quand un tour
de plus n'ajoute rien --- c'est la version abstraite de \textsc{CheckNoEffect}
partout. La \cref{tab:react-modele-correspondance} fait la correspondance
règle par règle.

\begin{table}[htbp]\centering\small
\begin{tabularx}{\linewidth}{@{}>{\raggedright}p{0.3\linewidth}>{\raggedright\arraybackslash}X@{}}
  \toprule
  React-tRace & \reactant{} (commit \code{e67b10a}) \\
  \midrule
  \textsc{SttBind} (phase \textsf{Init}) & amorçage hors boucle, \src{src/engine/fixpoint.rs}{314}{353} \\
  \textsc{AppSetComp}, \textsc{AppSetNormal} (mise en file) & \code{exec_setter_call} : join immédiat, \src{src/domains/interp/interpreter.rs}{341}{366} ; setter d'un autre composant : \code{SharedStateStore}, l.~368--392 \\
  \textsc{SttReBind} (rejeu de la file, $v_n \not\equiv v_0$) & lecture \code{StateVal} (store joint, référence \code{Versioned}), \src{src/domains/transfer/state_value.rs}{122}{134} ; « a changé ? » = test $N \lleq S$, \srcline{src/engine/fixpoint.rs}{495} \\
  \textsc{Eff}, \textsc{CommitEffsPath} & passes d'effets, tous les effets, \src{src/engine/fixpoint.rs}{415}{449} \\
  \textsc{StepCheck}, \textsc{CheckNoEffect} & test de convergence, \srcline{src/engine/fixpoint.rs}{495} \\
  cycle Check--Effect infini (§3.1.1) & \code{widen_trace}, \src{src/engine/fixpoint.rs}{514}{527}, et \code{effect_setter_writes}, l.~565--593 ; churn : lignes \code{exact} de \relation{effect_triggers} \\
  \textsc{EvalMult} infini (§3.1.2) & \regle{setter-in-render}, \Error{} si l'appel domine toutes les sorties \\
  \textsc{StepEvent} & \code{HookEntry::Handler}, passes de handlers hors \code{widen_trace}, \src{src/engine/fixpoint.rs}{451}{483} \\
  setter $\langle \ell, p\rangle$ & \code{Expr::StateSetter} s'évalue en \code{component_setter(composant, label)}, \srcline{src/domains/transfer/state_value.rs}{135} (ADR-012 : \og Corresponds to \code{Set_clos { label, path }} of React-tRace\fg{}) \\
  \textsc{AppCom}, \textit{init}, \textit{reconcile} & \code{eval_comp_app} : analyse descendante des enfants avec leurs props abstraites, \src{src/domains/transfer/state_value.rs}{470}{594} \\
  \bottomrule
\end{tabularx}
\caption{Table de correspondance entre les règles de React-tRace et les mécanismes de \reactant.}
\label{tab:react-modele-correspondance}
\end{table}

\subsection{Ce qu'on garde, ce qu'on abstrait}

\begin{table}[htbp]\centering\small
\begin{tabularx}{\linewidth}{@{}>{\raggedright}p{0.28\linewidth}>{\raggedright}X>{\raggedright\arraybackslash}X@{}}
  \toprule
  Fait de React & Ce que \reactant{} en fait & Pourquoi c'est sûr \\
  \midrule
  Le rendu, les effets et les handlers s'exécutent à des moments distincts & gardé : un CFG par corps, trois sortes de passes & structure \\
  Un effet capture l'environnement de son rendu & gardé : chaque effet part de l'environnement de sortie du rendu & exact \\
  L'identité d'un hook est sa position & gardé : label = ordre d'appel & exact \\
  Un setter est une valeur liée à un slot d'un composant & gardé : \code{component_setter}, \code{SharedStateStore} & exact \\
  File d'updaters, ordre, batching & abstraits : join immédiat, updater sur le store & le join contient tout résultat de file \\
  \og l'effet tourne ssi une dep a changé\fg{} & abstrait : tous les effets à chaque itération ; deps relues par les règles & plus d'exécutions que React \\
  Bail-out \code{Object.is} & abstrait : tout set peut re-rendre & plus de rendus que React \\
  Ordre enfants-parent des effets & abstrait : join commutatif, effets croisés vus à l'itération suivante & ordre quelconque couvert \\
  Nombre de clics & abstrait : handlers exécutés 0..N fois dans le point fixe & toute séquence couverte \\
  DOM, réconciliation fine, Strict Mode & ignorés (un élément natif vaut \code{Stable}) & sans effet sur l'état \\
  Durée d'une suite de rendus & non calculée & voir le piège de \code{Toggle} \\
  \bottomrule
\end{tabularx}
\caption{Les faits de React et leur sort dans l'analyseur.}
\label{tab:react-modele-garde-abstrait}
\end{table}

La \cref{tab:react-modele-garde-abstrait} résume les choix. L'argument de
soundness qui les réunit tient en une phrase : chaque abstraction
\emph{ajoute} des comportements (plus d'effets exécutés, plus de rendus, plus
de valeurs dans chaque slot), jamais n'en retire, et la mise à jour faible
garantit que ce qu'un tour concret peut écrire est contenu dans le store. La
définition précise de la soundness est au \cref{chap:interpretation-abstraite}
(\cref{def:soundness}) ; la dernière ligne du tableau en marque la
portée : ce qui est sur-approximé, c'est l'ensemble des états, pas la
terminaison.

\subsection{Les callbacks : dans le cycle ou hors du cycle}

Un effet qui appelle un setter dans un \code{.then} ou un \code{setTimeout}
relance un tour sans intervention extérieure ; un effet qui s'abonne avec
\code{addEventListener} ne le fait qu'à chaque événement. L'analyseur doit
distinguer les deux, sinon tout compteur de clics serait une boucle infinie.

% Table literate locale : « … » (l. 13) casse la table commune (\dots) sous pdflatex.
\rustsnippet[literate={…}{{...}}1 {→}{{$\rightarrow$}}1]{src/domains/interp/callbacks.rs}{6}{23}{la classe de déclenchement d'un callee}

La classe est décidée par le nom du callee (\code{classify_callee},
\src{src/domains/interp/callbacks.rs}{28}{56}) : \code{setTimeout},
\code{setInterval}, \code{queueMicrotask}, \code{requestAnimationFrame},
\code{then}/\code{catch}/\code{finally} et les fonctions d'ordre supérieur
synchrones (\code{map}, \code{forEach}\dots) sont \code{InCycle} et leurs
arguments fonctionnels sont exécutés ; \code{addEventListener} est
\code{Subscription} ; tout autre callee est \code{Unknown}.

\begin{decision}[ADR-009 --- points d'entrée et classe de déclenchement]
\emph{Décision} (\adr{9}). La descente dans les callbacks est
\emph{sémantique} (le store bouge), pas seulement structurelle ; les handlers
JSX et les abonnements sont des \emph{points d'entrée} séparés, exécutés 0..N
fois dans le point fixe, dont les écritures entrent dans le store mais pas dans
\code{widen_trace} : « clicking 1000$\times$ isn't a bug ». \emph{Alternative
refusée.} Descendre uniformément dans tout callback aurait armé la valeur sur
\code{addEventListener('click', () => setCount(c => c + 1))} dans un effet, un
faux positif d'\regle{infinite-loop} sur du React correct. \emph{Choix
explicite} : \code{Unknown} n'est pas descendu, \og for a linter, false
positives are more costly than false negatives [\dots] We accept the FN\fg{}.
\end{decision}

\begin{piege}
Le choix \og \code{Unknown} $\to$ skip\fg{} de l'ADR-009 est en tension avec
l'invariant du projet, « faux négatifs INTERDITS ». Un setter appelé dans
\code{myUtil(() => setX(x + 1))} ne fait pas bouger le store ; un mécanisme
ultérieur (l'inlining des fonctions locales connues) réduit ce trou sans le
fermer. La relation des écrivains, elle, garde ces écritures avec une phase
$\Top$ (\code{SetterCallPhase::Unknown}), que \regle{setter-in-render} signale
en \Warning. Le store et la relation n'ont donc pas la même politique : « absent
du store » ne veut pas dire « jamais écrit ».
\end{piege}

Le moment d'une écriture, vu par la relation des écrivains
(\cref{chap:relations}), est une phase :

\rustsnippet{src/engine/setters.rs}{57}{70}{les phases d'un appel de setter}

\code{Handler} dit qu'un événement utilisateur s'interpose ; \code{Deferred},
qu'un tour ultérieur de la boucle d'événements s'interpose (minuterie,
microtâche, cleanup). Depuis l'ADR-035 (\adr{35}), une écriture placée après un
\code{await} est aussi \code{Deferred} : le lowering coupe le bloc à l'\code{await}
par une arête \code{EdgeKind::Await}, et la marche des setters promeut
\code{Sync} en \code{Deferred} dans les blocs qui suivent. C'est la traduction
exacte d'un fait du langage : après un \code{await}, on est sorti de toute
phase de React.

\subsection{Trois écarts vérifiés}

L'exactitude impose de montrer aussi où la sémantique abstraite du commit
étudié \emph{ne} sur-approxime \emph{pas} la sémantique concrète, ou la
sur-approxime au mauvais niveau. Ce sont des défauts, pas des choix ; ils sont
présentés ici parce que chacun se comprend par un fait React de ce chapitre.

\begin{piege}
\textbf{L'updater fonctionnel sur le chemin programme (faux négatif).}
\begin{lstlisting}[language=TypeScript,caption={\code{D} : un updater dans un effet gardé par l'état qu'il écrit (\file{/tmp/ch02/updater.tsx}, imports omis)},label={lst:react-modele-updater}]
export function D() {
  const [count, setCount] = useState(0);
  useEffect(() => {
    setCount((c) => c + 1);
  }, [count]);
  return <div>{count}</div>;
}
\end{lstlisting}
En React, c'est une boucle certaine. La CLI ne dit rien :
\begin{sortie}
  [verbose] D: 2 iteration(s), widened: []
✓  1 file(s) no issues found.
\end{sortie}
Pourtant le test \code{functional_updater_with_state_dep_is_flagged} de
\file{tests/functional_updater.rs}, qui analyse le même composant, passe
(\code{cargo test --test functional_updater} : \code{5 passed}). Le test appelle
\code{analyze_component} (analyse intra, \code{inter = None}) ; la CLI passe par
l'analyse de programme (\code{inter} présent). La seconde branche
d'\code{exec_setter_call} (\src{src/domains/interp/interpreter.rs}{368}{392})
s'applique alors à tout callee qui s'évalue en setter, y compris le setter du
composant \emph{lui-même}, et joint dans le \code{SharedStateStore}
l'argument \emph{évalué}, c'est-à-dire la fermeture \code{c => c + 1}
(\code{PerRender}) et non la valeur qu'elle renvoie ; ce store est rejoint au
test de convergence, le slot devient un mélange nombre--référence sur lequel
\code{c + 1} vaut $\Top$, et plus rien ne croît.
\end{piege}

\begin{piege}
\textbf{\code{useReducer} (faux positif au niveau \Error).} Parce que le
réducteur est jeté (\cref{sec:react-modele-usestate}), \code{dispatch(action)}
écrit l'action dans le slot :
\begin{lstlisting}[language=TypeScript,caption={\code{R} : un réducteur qui ignore l'action \code{noop} (\file{/tmp/ch02/reducer.tsx}, imports omis)},label={lst:react-modele-reducer}]
function reducer(s: { n: number }, action: { type: string }) {
  if (action.type === "noop") return s;
  return { n: s.n + 1 };
}

export function R() {
  const [state, dispatch] = useReducer(reducer, { n: 0 });
  useEffect(() => {
    dispatch({ type: "noop" });
  }, [state]);
  return <div>{state.n}</div>;
}
\end{lstlisting}
\begin{sortie}
  R  (2 hooks)  reducer.tsx
    error  infinite-loop  [hook:0]  (line 10:2)  this effect recreates object state `state` it depends on. Every run stores a fresh reference (`Object.is` always fails) and re-triggers itself: infinite render loop
       → a fresh value is written to state `state` here [hook:1] (line 11:4)
\end{sortie}
En React, le réducteur renvoie \code{s} inchangé, \code{Object.is} conclut à
l'égalité, pas de rendu : aucune boucle. L'analyseur voit un objet frais (l'action)
écrit dans le slot dont l'effet dépend, et conclut au churn. C'est une \Error{}
fausse, ce que \file{docs/limitations.md} exclut (\og A false positive never
carries an Error\fg{}).
\end{piege}

\begin{piege}
\textbf{La boucle oscillante (faux négatif).} C'est \code{Toggle}
(\cref{ex:canon-toggle}) : la valeur abstraite converge parce que
l'ensemble des états est fini, et aucune preuve structurelle ne couvre une
négation booléenne.
\end{piege}

\section{Les exemples canoniques du manuscrit}
\label{sec:react-modele-exemples}

Les chapitres suivants reviennent sans cesse à une poignée de programmes. Ils
sont définis dans ce chapitre (ou, pour \code{Timer}, dans l'introduction) et
réunis dans la \cref{tab:react-modele-canon}, avec le comportement de React et
le verdict de \reactant{} au commit \code{e67b10a}. Tous les fichiers sont sous
\file{/tmp/ch02/} et les sorties ont été obtenues par
\code{reactant check --no-color --info --show-clean --project plain <fichier>}.

Il manquait au chapitre la sortie de \code{Timer}. Le fichier
\file{/tmp/ch02/timer.tsx} contient une variante du \code{Timer} de l'introduction
(\cref{lst:introduction-timer}) : même structure, un état nommé \code{n} et un
callback écrit en une ligne, \code{setInterval(() => setN(n + 1), 1000)} ; et sa
correction
\code{TimerFixed}, qui passe l'updater \code{(m) => m + 1} :

\begin{sortie}
  Timer  (2 hooks)  timer.tsx
    warn   missing-deps  [hook:1]  var:n  (line 5:2)  `n` is used in this effect but not in its deps array, and it is recreated on every render
    error  stale-closure  [hook:1]  var:n  (line 6:10)  the `setInterval` callback registered by this mount-only effect reads `n` and writes it back. `n` was captured once at mount, so every firing recomputes from the same frozen value and the state can never advance past its first update
    info   widening-info  (line 5:2)  state `n` kept changing during analysis and was approximated to converge, so findings that depend on it may be imprecise
  TimerFixed  (2 hooks)  timer.tsx  ✓
    verified  stale-closure  no long-lived callback captures a stale state value
\end{sortie}

Le callback de \code{setInterval} est \code{InCycle}
(\cref{sec:react-modele-concrete}) : son écriture entre dans le point fixe, d'où
l'élargissement. Le fait React qui fonde l'\Error{} est celui de la
\cref{sec:react-modele-usestate} : \code{n} est une constante du rendu de
montage, et l'effet, de deps \code{[]}, ne tourne qu'une fois.

\begin{center}\small
\begin{longtable}{@{}>{\raggedright}p{0.25\linewidth}>{\raggedright}p{0.3\linewidth}>{\raggedright\arraybackslash}p{0.37\linewidth}@{}}
  \caption{Les exemples canoniques. « Ex. » renvoie à l'encadré ou à l'extrait qui les définit.}\label{tab:react-modele-canon}\\
  \toprule
  Composant (fichier) & À l'exécution (React) & \reactant{} (\code{e67b10a}) \\
  \midrule
  \endfirsthead
  \toprule
  Composant (fichier) & À l'exécution (React) & \reactant{} (\code{e67b10a}) \\
  \midrule
  \endhead
  \code{Counter} (\file{counter.tsx}), \cref{ex:canon-counter} & compte les clics & aucun diagnostic \\
  \code{Batch} (\file{batching.tsx}), \cref{lst:react-modele-batch} & \code{+2} par clic & aucun diagnostic \\
  \code{SameValue} (\file{samevalue.tsx}), \cref{lst:react-modele-samevalue} & bail-out, un seul rendu & \Warning{} \regle{redundant-set-state} \\
  \code{Todos} (\file{lazy.tsx}), \cref{lst:react-modele-lazy} & appel inutile à chaque rendu & \Warning{} \regle{lazy-init} \\
  \code{EveryCommit}, \code{OnChange} (\file{deps.tsx}) & boucle infinie & \Warning{} \regle{infinite-loop} \\
  \code{MountOnly} (\file{deps.tsx}) & deux rendus puis arrêt & \Warning{} \regle{missing-deps}, \Info{} \regle{widening-info} \\
  \code{Inf} (\file{selfcounter.tsx}), \cref{ex:canon-selfcounter} & boucle infinie & \Warning{} \regle{infinite-loop} \\
  \code{SelfCounter} (id.) & quatre rendus puis arrêt & \Info{} \regle{widening-info} seulement \\
  \code{Flicker} (\file{flicker.tsx}), \cref{ex:canon-flicker} & un rendu de trop au montage & \Warning{} \regle{unnecessary-rerender} \\
  \code{Toggle} (\file{toggle.tsx}), \cref{ex:canon-toggle} & boucle infinie (oscillation) & \textbf{aucun diagnostic} (FN) \\
  \code{Width} (\file{listener.tsx}), \cref{lst:react-modele-listener} & abonnements empilés & \Warning{} \regle{missing-cleanup} \\
  \code{Search} (\file{search.tsx}), \cref{lst:react-modele-search} & un effet repart à chaque rendu & \Warning{} \regle{always-unstable-deps} \\
  \code{ObjChurn} (\file{objchurn.tsx}), \cref{ex:canon-objchurn} & boucle infinie (identité) & \Error{} \regle{infinite-loop} \\
  \code{ObjStateDep} (id.) & effet de montage seul & aucun diagnostic \\
  \code{Cond} (\file{cond.tsx}), \cref{ex:canon-cond} & interface cassée après le clic & \Error{} \regle{conditional-hook} \\
  \code{Inf2} (\file{setter_in_render.tsx}), \cref{ex:canon-inf2} & \og Too many re-renders\fg{} & \Error{} \regle{setter-in-render} \\
  \code{List} (\file{adjust.tsx}), \cref{lst:react-modele-adjust} & motif documenté, converge & \Warning{} \regle{setter-in-render} sur \code{setSelection} (FP) \\
  \code{Impure} (\file{strict.tsx}), \cref{lst:react-modele-impure} & rendu impur (Strict Mode le révèle) & aucun diagnostic \\
  \code{ThemeProvider}, \code{Button} (\file{context.tsx}) & consommateurs re-rendus & \Warning{} \regle{unstable-context-value}, \Info{} \regle{analysis-limit} \\
  \code{Page} (\file{app/page.tsx} de \file{/tmp/ch02next}) & hook dans un Server Component & \Warning{} \regle{server-component-hook} \\
  \code{Timer} (\file{timer.tsx}), \cref{lst:introduction-timer} & chronomètre bloqué à 1 & \Error{} \regle{stale-closure}, \Warning{} \regle{missing-deps} \\
  \code{TimerFixed} (id.) & correct & aucun diagnostic \\
  \code{D} (\file{updater.tsx}), \cref{lst:react-modele-updater} & boucle infinie & \textbf{aucun diagnostic} (FN) \\
  \code{R} (\file{reducer.tsx}), \cref{lst:react-modele-reducer} & pas de boucle (bail-out) & \textbf{\Error{}} \regle{infinite-loop} (FP) \\
  \bottomrule
\end{longtable}
\end{center}

\begin{resume}
\begin{itemize}
  \item Un composant est une fonction que React appelle autant de fois qu'il
    le veut. Un tour : déclencheur, rendu, commit, effets ; un set dans un
    effet relance un tour, un set dans un handler attend un événement, un set
    pendant le rendu fait jeter le rendu (\cref{fig:react-modele-chronologie}).
  \item \code{set} planifie le prochain rendu ; les mises à jour sont mises en
    file et regroupées ; un updater lit la file ; \code{Object.is} fait sauter
    un rendu inutile. L'analyseur ne garde ni file ni ordre : une écriture est
    un join dans le \code{StateStore}, un updater est exécuté sur la valeur du
    store (\code{exec_setter_call}).
  \item Un effet tourne après le commit si une dep a changé selon
    \code{Object.is}. Le point fixe exécute \emph{tous} les effets à chaque
    itération ; les deps sont relues par les règles et par la relation
    \relation{effect_triggers} (bit \code{exact} : l'effet \emph{doit} repartir).
  \item L'identité compte autant que la valeur : littéraux \code{PerRender},
    lecture d'un état \code{Versioned} par son slot (double vue de l'état,
    ADR-017), mémos comme join des deps. \code{ObjChurn} est une \Error{} par
    trois faits must, pas par divergence.
  \item Un hook est identifié par sa position (\code{HookLabel}) ; un hook
    conditionnel ne domine pas toutes les sorties ; la fin d'un corps est un
    \code{return} (ADR-025).
  \item Strict Mode n'est pas modélisé, le contexte côté consommateur est
    $\Top$, les Server Components sont analysés quand même et un hook y est un
    \Warning{} (ADR-026).
  \item React-tRace fournit la sémantique concrète de référence (ADR-001) :
    mémoire d'arbre, décisions \textsf{Check}/\textsf{Effect}, quatre
    transitions \textsc{StepInit}, \textsc{StepEffect}, \textsc{StepCheck},
    \textsc{StepEvent} ; la \cref{tab:react-modele-correspondance} la relie au
    code. Le contrat est conceptuel : ni \file{docs/semantics.md}, ni oracle de
    tests.
  \item Le point fixe sur-approxime l'ensemble des états, pas la durée d'une
    suite de rendus (\code{Toggle}) ; trois écarts vérifiés au commit étudié :
    \code{Toggle} et \code{D} (FN), \code{R} (\Error{} fausse).
\end{itemize}
{\raggedright Fichiers rencontrés : \file{src/engine/fixpoint.rs},
\file{src/domains/interp/interpreter.rs}, \file{src/domains/interp/callbacks.rs},
\file{src/domains/transfer/state_value.rs}, \file{src/domains/impls/stability.rs},
\file{src/domains/stores/state_store.rs}, \file{src/engine/triggers.rs},
\file{src/engine/setters.rs}, \file{src/ir/hooks.rs},
\file{src/lowering/hook_extractor.rs}. Le \cref{chap:interpretation-abstraite}
donne le cadre mathématique (treillis, concrétisation, point fixe, widening,
soundness) dans lequel toutes ces approximations se formulent exactement.\par}
\end{resume}

\section*{Exercices}
\addcontentsline{toc}{section}{Exercices}

\begin{exercice}{Les itérations d'\code{Inf}}{react-modele-inf}
Pour \code{Inf} (\cref{ex:canon-selfcounter}), donner la valeur du slot
\code{s} (composante numérique) au début de chaque itération de
l'\cref{alg:react-modele-fixpoint}, sachant que \code{widen_threshold = 3}, et
dire à quelle itération \code{widen_trace} reçoit le slot. Comparer avec la
sortie \code{--verbose} : \code{Inf: 3 iteration(s), widened: [0]}.
\end{exercice}
\begin{solution}
Amorçage : $[0,0]$. Passe 1 : l'effet écrit $[1,1]$, $N = [0,1] \not\lleq [0,0]$,
$i = 1$, $S \leftarrow [0,1]$. Passe 2 : écriture $[1,2]$, $N = [0,2]$, $i = 2$,
$S \leftarrow [0,2]$. Passe 3 : écriture $[1,3]$, $N = [0,3]$, $i = 3 \geq 3$ :
le slot entre dans \code{widen_trace} avec l'itération 3 (d'où « widened at
iteration 3 ») et $S \leftarrow [0,2] \widen [0,3]$, non borné à droite. Passe
4 : plus rien ne croît, sortie. Le compteur vaut 3.
\end{solution}

\begin{exercice}{Deux updaters}{react-modele-deux-updaters}
Un handler exécute \code{setN((c) => c + 1); setN((c) => c + 1);} avec un
store $\{n \mapsto [0,0]\}$. Que vaut le slot après la passe selon
\code{exec_setter_call} ? Que donne React ? Pourquoi le résultat abstrait
est-il correct bien qu'il contienne une valeur que React ne produit jamais
après ce clic ?
\end{exercice}
\begin{solution}
Premier updater : $c = [0,0]$, écriture $[1,1]$, slot $[0,1]$. Second :
$c = [0,1]$, écriture $[1,2]$, slot $[0,2]$. React rejoue la file :
$0 \to 1 \to 2$, soit \code{2}. $[0,2]$ contient \code{2} ; les valeurs \code{0}
et \code{1} en trop sont le prix du join (imprécision), pas une erreur : une
sur-approximation peut contenir des états impossibles, elle ne doit pas en
omettre de possibles.
\end{solution}

\begin{exercice}{Tous les effets}{react-modele-tous-effets}
En vous appuyant sur les règles \textsc{CommitEffsPath} et
\textsc{CommitEffsPathIdle}, montrer qu'exécuter tous les effets à chaque
itération (\cref{lst:react-modele-effect-pass}) produit un store qui contient
toutes les écritures des effets que React exécute. Qu'est-ce que cette
sur-approximation coûte en précision, et où le projet la récupère-t-il ?
\end{exercice}
\begin{solution}
React exécute un sous-ensemble des effets (ceux dont la vue porte
\textsf{Effect}) ; l'analyseur les exécute tous, sur un store qui contient déjà
toutes les valeurs atteignables à ce point ; comme les écritures sont jointes,
l'ensemble obtenu contient celui de n'importe quel sous-ensemble. Le coût : un
effet de montage semble tourner à chaque tour (\code{MountOnly} est élargi). Le
projet récupère la précision dans les règles, qui relisent les deps
(\regle{infinite-loop} écarte \code{[]} et les deps toutes stables) et dans
\relation{effect_triggers}.
\end{solution}

\begin{exercice}{Les trois faits de l'\Error{} d'\code{ObjChurn}}{react-modele-triple-must}
Énumérer les trois faits must qui fondent l'\Error{} d'\code{ObjChurn}. Si la
dep est \code{[obj.a]} au lieu de \code{[obj]}, lequel manque, et pourquoi ?
\end{exercice}
\begin{solution}
Must-change : la dep est le slot exact et la valeur écrite est \code{PerRender}.
Must-reach : l'écriture a lieu sur tous les chemins du corps de l'effet.
Must-rerun : React ré-exécute un effet dont une dep a changé. Avec
\code{[obj.a]}, le premier manque : \code{{...obj, b: 2}} conserve la valeur
du champ \code{a}, la dep peut ne pas changer ; l'ADR-017 plafonne ce cas au
\Warning.
\end{solution}

\begin{exercice}{\code{ObjStateDep}}{react-modele-objstatedep}
L'état de \code{ObjStateDep} est initialisé par un littéral objet, donc le store
contient une référence \code{PerRender}. Expliquer pourquoi
\regle{always-unstable-deps} ne tire pas, en citant la ligne de code
responsable, et dire quelle hypothèse cette ligne fait sur le programme.
\end{exercice}
\begin{solution}
La lecture \code{Expr::StateVal} remplace la composante référence par
\code{Versioned({(composant, label)})}
(\cref{lst:react-modele-stateval}, l.~130--132) ; la règle exige
\code{PerRender}. L'hypothèse : les sets ont lieu hors du rendu ; sa violation
est signalée par \regle{setter-in-render} (soundness en couches).
\end{solution}

\begin{exercice}{Pourquoi \code{Toggle} converge}{react-modele-toggle}
Calculer la suite des valeurs booléennes abstraites du slot \code{on} de
\code{Toggle} et expliquer la sortie \code{1 iteration(s), widened: []}. Quel
fait \emph{structurel}, à la manière du bras churn, permettrait de prouver la
boucle ? (Question ouverte : il n'existe pas de réponse dans le code.)
\end{exercice}
\begin{solution}
Amorçage \code{false} ; la passe d'effet écrit \code{!on = true} ; le slot
devient \code{{false, true}}, c'est-à-dire le top du treillis booléen ;
$i = 1$ ; la passe suivante n'ajoute rien, sortie. Aucune croissance, donc pas
de widening. Il faudrait établir que la dep est le slot exact (ligne
\code{exact} de \relation{effect_triggers}) et que la valeur écrite diffère
\emph{toujours} de la valeur lue (ici $\lnot v \neq v$ pour tout booléen), sur
tous les chemins de l'effet : un must-change par valeur plutôt que par identité.
\end{solution}

\begin{exercice}{\code{Inf2} de bout en bout}{react-modele-inf2}
Pour \code{Inf2} (\cref{ex:canon-inf2}), donner la règle de React-tRace
qui s'applique à l'appel \code{setS((x) => x)}, le théorème qui explique la
boucle, et le fait de CFG que \regle{setter-in-render} utilise pour émettre une
\Error{} plutôt qu'un \Warning.
\end{exercice}
\begin{solution}
\textsc{AppSetComp} (set pendant l'évaluation du corps) pose \textsf{Check} ;
par le \cref{thm:trace-reevaluation}, l'évaluation est retentée
(\textsc{EvalMult}), et comme chaque essai rappelle le setter, indéfiniment,
quelle que soit la valeur. \regle{setter-in-render} émet une \Error{} parce
que le bloc de l'appel domine toutes les sorties du rendu : l'appel a lieu à
chaque rendu.
\end{solution}
